% !TeX root = ../../Book3.tex
% 纲要标题：## 第14章　椭圆边值问题与参数族
% 纲要位置：几何学计划纲要.md，第 2669 行。
% 正文状态：初稿；尚未正式审核。

% 纲要内容开始
% 
% > 本章目标与证明范围：把原来一节里的边界、参数和系统拆为有明确对象的整章。主证明限定光滑紧区域上的 Dirichlet、Neumann 及所选丛值问题；一般伪微分边界理论只作准确出口。
% 
% 纲要内容结束

\chapter{椭圆边值问题与参数族}
\label{g3:ch:14}
\BookChapterNumber{g3:ch:14}

\begin{chapterabstract}
	在有边界区域上，微分表达式与边界条件共同确定算子。边界既改变弱方程的测试空间，也决定哪些整体估计能够成立。本章由 Green 配对建立 Dirichlet、Neumann 与混合条件的能量解，随后用边界展平、切向差商和法向导数恢复证明光滑边界上的正则性，并构造非齐次数据所需的迹与提升。半空间衰减模型说明补条件怎样约束边界数据；所选丛值问题以 Hodge 的绝对、相对条件为核心，落实 Fredholm 可解性。最后讨论参数变化中的投影、Green 算子与统一估计，区分固定核维数的平稳求解和穿过共振值时的核跳跃。
\end{chapterabstract}

\begin{learninggoals}
	\begin{itemize}
		\item 从 Green 公式辨认外共法向通量、弱测试空间及 Neumann 兼容条件的符号。
		\item 构造非齐次 Dirichlet 提升，证明边界弱解的存在、唯一性和能量刻画。
		\item 用切向差商取得贴边导数，再从方程恢复法向导数，完成高阶边界估计。
		\item 说明半整数阶迹空间和边界数据提升为何与二阶正则性配套。
		\item 计算半空间的衰减解及边界符号，检验补条件并处理绝对、相对形式边界条件。
		\item 在固定空间中构造参数核投影与 Green 算子，核对统一常数、兼容条件与边界域识别。
		\item 用角点、条件交界与共振例子辨明整体正则性和参数连续性的适用范围。
	\end{itemize}
\end{learninggoals}

在区间 \([0,1]\) 上，\(-u''=0\) 的所有解都是 \(u=ax+b\)。加入 \(u(0)=u(1)=0\)，只剩零解；改成 \(u'(0)=u'(1)=0\)，则留下全部常数。若要求 \(-u''=1\) 且两端导数为零，积分给出 \(0=1\)，因而无解。同一个内部表达式既可能唯一可解，也可能留有核或遇到兼容障碍，区别来自边界。

\ProseParagraphBreak

边界还影响正则性证明。内部差商可以沿每个方向平移，贴近边界时，法向平移可能离开区域，或破坏指定的边界值。将边界展平以后，切向平移保留半空间；先取得这些方向的导数，再借助方程中非退化的法向二阶项恢复剩余导数，便形成可继续迭代的证明。若边界有角，或者 Dirichlet 与 Neumann 条件在一点相接，平移所保留的结构发生变化，显式分数次幂解会显示相应的正则性损失。

\ProseParagraphBreak

参数问题则要求同时跟踪算子与其定义域。例如在零端点条件下，\(-u''-t u\) 作用于 \(\sin(\pi x)\) 的结果为 \((\pi^2-t)\sin(\pi x)\)。当 \(t\) 接近 \(\pi^2\) 时，这个方向上的逆越来越大；到达 \(\pi^2\) 时，它转化为核方向，并产生一个新的右端正交条件。单个参数处的椭圆估计，必须配合核、投影及边界条件的统一控制，才能用于变化的几何族。

\ProseParagraphBreak

所需基础是第十一章的迹、紧嵌入和弱紧性，第十二章的能量存在性及内部正则性，以及第十三章的 Fredholm 分解、紧自伴谱与 Banach 隐函数方法。本章先完成标量边界问题，再证明所选形式丛问题；最后一节以固定 Banach 空间中的有限维约化组织参数依赖。基本变分与边界正则性可对照 Brezis 第9.5—9.6节\cite[第9.5—9.6节]{Brezis2011}；一般边界伪微分方法在所选证明之后说明其所处理的对象。

% !TeX root = ../../Book3.tex
% 纲要标题：### 14.1 边值问题的弱形式
% 纲要位置：几何学计划纲要.md，第 2673 行。
% 正文状态：初稿；尚未正式审核。

% 纲要内容开始
% 
% * Green 公式和边界配对
% * Dirichlet 变分问题
% * Neumann 兼容条件及常数核
% * 混合边界条件与角点奇性
% 
% 纲要内容结束

\section{边值问题的弱形式}
\label{g3:sec:1401}

边界条件决定了方程允许怎样变化。指定函数值时，变分在边界消失；指定通量时，边界上的测试仍然自由，通量便成为右端的一部分。本节从这一差别构造弱问题，先取得能量空间中的解，再把边界处导数的恢复留给下一节。

\subsection{Green 公式和边界配对}
\label{g3:subsec:1401:01}
以下以紧的光滑带边界黎曼流形 \((M,g)\) 为区域，\(\dim M=n\ge1\)，每个连通分支都与 \(\partial M\) 相交。Euclid 空间中光滑有界域的闭包是基本例子。记外单位法向为 \(\nu\)，体积与边界密度分别为 \(\mu_g,\mu_{\partial}\)。一维时边界密度就是端点计数测度。系数 \(A:\mathrm TM\to\mathrm TM\) 光滑、实对称，并满足
\[
	\lambda|\xi|^2\le\langle A\xi,\xi\rangle_g\le\Lambda|\xi|^2,
	\qquad0<\lambda\le\Lambda<\infty.
\]
它按复线性延拓作用于复梯度。对实光滑函数 \(c\)，取
\begin{equation}\label{g3:eq:boundary-operator-form}
	\begin{aligned}
		Pu&=-\operatorname{div}_g(A\operatorname{grad}_g u)+cu,\\
		Q(u,v)&=\int_M\bigl(\langle A\operatorname{grad}u,\operatorname{grad}v\rangle_g
		+cu\overline v\bigr)\,\mu_g,\\
		B_Au&=\langle A\operatorname{grad}u,\nu\rangle_g.
	\end{aligned}
\end{equation}
内积对第一变量线性。\(B_Au\) 称为外\term[gongfaxiang daoshu]{共法向导数}[conormal derivative]；只有 \(A=I\) 时它就是普通外法向导数。

\begin{proposition}\label{g3:prop:boundary-green-identities}
	设 \(M\) 为紧光滑带边界黎曼流形，\(A\) 为光滑实对称丛自同态，\(c\) 为实光滑函数，令 \(\gamma\) 为边界迹，\(\nu\) 为外单位法向，\(P=-\operatorname{div}(A\operatorname{grad})+c\)，\(Q(u,v)=\int_M(\langle A\operatorname{grad}u,\operatorname{grad}v\rangle+cu\overline v)\,\mu_g\)，\(B_Au=\langle A\operatorname{grad}u,\nu\rangle\)，其中 \(\mu_g\) 和 \(\mu_\partial\) 分别为体积密度及边界密度。若 \(u,v\in C^\infty(M;\mathbb C)\)，则
	\begin{align}
		(Pu,v)_2&=Q(u,v)-\int_{\partial M}B_Au\,\overline{\gamma v}\,\mu_{\partial},
		\label{g3:eq:boundary-green-first}\\
		(Pu,v)_2-(u,Pv)_2
		&=\int_{\partial M}\bigl(\gamma u\,\overline{B_Av}
		-B_Au\,\overline{\gamma v}\bigr)\,\mu_{\partial}.
		\label{g3:eq:boundary-green-second}
	\end{align}
	这里 \(\gamma\) 表示边界限制。
\end{proposition}
\begin{proof}
	对复向量场 \(\overline v A\operatorname{grad}u\) 逐实、虚部使用散度公式。乘积法则给出
	\[
		\operatorname{div}(\overline v A\operatorname{grad}u)
		=\overline v\operatorname{div}(A\operatorname{grad}u)
		+\langle A\operatorname{grad}u,\operatorname{grad}v\rangle,
	\]
	边界通量为 \(B_Au\,\overline{\gamma v}\)。移项并加入零阶项就得到第一式。交换 \(u,v\) 后取共轭，利用 \(Q(u,v)=\overline{Q(v,u)}\)，两式相减即得第二式。这里使用密度的散度公式，不要求 \(M\) 可定向。
\end{proof}

式\eqref{g3:eq:boundary-green-second}把函数值与通量成对放在边界上；检验边界实现是否对称，就要检查这个配对是否消失。对 \(H^1\) 函数，第十一章的\cref{g3:prop:sobolev-boundary-bundles}给出有界迹 \(\gamma:H^1(M)\to L^2(\partial M)\)，且核为 \(H_0^1(M)\)。用迹的实际像记录可指定的边界值：
\begin{definition}\label{g3:def:trace-range-space}
	设 \(M\) 为紧光滑带边界黎曼流形，\(\gamma:H^1(M)\to L^2(\partial M)\) 为边界迹。定义
	\[
		\mathcal T(\partial M)=\gamma H^1(M),\qquad
		\|\psi\|_{\mathcal T}=\inf_{\gamma w=\psi}\|w\|_{H^1(M)}.
	\]
	记 \(\mathcal T^\dagger\) 为其连续反线性泛函空间；实值问题中取连续线性泛函。
\end{definition}
由于 \(H_0^1\) 是闭子空间，\(\mathcal T\) 等距同构于 Hilbert 商空间 \(H^1/H_0^1\)。具体地，把任意提升投影到 \((H_0^1)^{\perp_{H^1}}\)，就得到唯一的最小 \(H^1\) 范数提升。这个操作给出有界线性右逆 \(E:\mathcal T\to H^1\)，满足 \(\gamma E=I\)。因此 \(\mathcal T\) 完备，并连续嵌入边界 \(L^2\)。下一节将用 Fourier 方法把它识别为 \(H^{1/2}(\partial M)\)。
\symidx{\(\mathcal T(\partial M)\)：赋商范数的 \(H^1\) 边界迹空间}

\begin{proposition}\label{g3:prop:weak-conormal-trace}
	设 \(M\) 为紧光滑带边界黎曼流形，\(A\) 光滑、实对称且一致正定，\(c\) 实光滑。令 \(\gamma\) 为边界迹，\(\nu\) 为外单位法向，\(P=-\operatorname{div}(A\operatorname{grad})+c\)，\(Q(u,v)=\int_M(\langle A\operatorname{grad}u,\operatorname{grad}v\rangle+cu\overline v)\,\mu_g\)，\(B_Au=\langle A\operatorname{grad}u,\nu\rangle\)，其中 \(\mu_g\) 和 \(\mu_\partial\) 分别为体积密度及边界密度。记 \(\mathcal T=\gamma H^1(M)\) 配商范数 \(\inf_{\gamma w=\psi}\|w\|_{H^1}\)，\(\mathcal T^\dagger\) 为其有界反线性对偶。若 \(u\in H^1(M;\mathbb C)\)，且 \(Pu=f\in L^2(M)\) 在分布意义下成立，则有唯一 \(\gamma_Au\in\mathcal T^\dagger\) 满足
	\begin{equation}\label{g3:eq:weak-conormal-green}
		(\gamma_Au)(\gamma v)=Q(u,v)-(f,v)_2\qquad(v\in H^1(M)).
	\end{equation}
	并有 \(\|\gamma_Au\|_{\mathcal T^\dagger}\le C(\|u\|_{H^1}+\|f\|_2)\)。若 \(u\) 光滑，则该泛函为 \(\psi\mapsto\int_{\partial M}B_Au\,\overline\psi\,\mu_{\partial}\)。
\end{proposition}
\begin{proof}
	分布方程先给出 \(Q(u,v)=(f,v)_2\) 对内部紧支光滑 \(v\) 成立；两侧都对 \(H^1\) 范数连续，所以等式延拓到 \(H_0^1\)。因此右侧在 \(\ker\gamma\) 上为零，只依赖于 \(\gamma v\)，并且对所有提升的 \(H^1\) 范数有同一上界。取下确界即得商范数估计。\(\gamma\) 满射到 \(\mathcal T\)，所以泛函唯一。光滑情形由第一 Green 公式识别。
\end{proof}

这个定义没有对 \(H^1\) 函数的一阶导数直接取迹；它利用方程 \(Pu\in L^2\) 定义整体边界通量。Dirichlet 数据属于迹空间，而 Neumann 数据属于它的对偶，二者在弱问题中承担不同的配对。

\subsection{Dirichlet 变分问题}
\label{g3:subsec:1401:02}
在\term[Dirichlet bianjie tiaojian]{Dirichlet 边界条件}[Dirichlet boundary condition]下，指定 \(\gamma u=\psi\)，所有允许的变分都属于 \(H_0^1\)。\footnote{Johann Peter Gustav Lejeune Dirichlet（1805--1859），德国数学家，研究数论、Fourier 级数与位势论；指定边界值的调和函数问题以其命名。}先把所需的 Poincaré 估计写成带边界版本。

\begin{lemma}\label{g3:lem:boundary-poincare}
	设 \(M\) 为非空紧光滑带边界黎曼流形，维数至少一，每个连通分支均有非空边界。则有以下两种估计。
	\begin{enumerate}
		\item 若 \(u\in H^1(M)\) 在每个连通分支上积分为零，则 \(\|u\|_{H^1}\le C\|\mathrm du\|_2\)。
		\item 若可测集 \(\Gamma_D\subset\partial M\) 在每个分支的边界中均有正测度，且 \(\gamma u=0\) 于 \(\Gamma_D\)，则同样有 \(\|u\|_{H^1}\le C_D\|\mathrm du\|_2\)。特别地，该估计适用于 \(H_0^1(M)\)。
	\end{enumerate}
\end{lemma}
\begin{proof}
	两种条件分别定义 \(H^1\) 的闭子空间。若其中一种估计失败，可取满足该条件的 \(u_j\)，使 \(\|u_j\|_2=1\)、\(\|\mathrm du_j\|_2\to0\)。带边界的紧嵌入\cref{g3:prop:sobolev-boundary-bundles}与弱紧性给出子列，在 \(L^2\) 中强收敛、在 \(H^1\) 中弱收敛到 \(u\)。于是 \(\|u\|_2=1\)、\(\mathrm du=0\)，故 \(u\) 在各分支上为常数；这一结论由内部的\cref{g3:prop:weak-gradient-constant}得到，边界是零体积集。

	第一种情形中，各分支平均值在极限下仍为零，故全部常数为零。第二种情形中，有界线性迹保持弱收敛，限制到 \(\Gamma_D\) 的迹仍为零；各分支上边界集合的正测度迫使相应常数为零。两者都与 \(\|u\|_2=1\) 矛盾。加上一阶导数项即得到所列 \(H^1\) 估计。
\end{proof}

\begin{theorem}\label{g3:thm:weak-general-dirichlet}
	设 \(M\) 为非空紧光滑带边界黎曼流形，\(\dim M\ge1\)，每个分支均有边界；\(A\) 光滑、实对称且一致正定，\(c\in C^\infty(M;\mathbb R)\)、\(c\ge0\)。令 \(\gamma\) 为边界迹，\(\nu\) 为外单位法向，\(P=-\operatorname{div}(A\operatorname{grad})+c\)，\(Q(u,v)=\int_M(\langle A\operatorname{grad}u,\operatorname{grad}v\rangle+cu\overline v)\,\mu_g\)，\(B_Au=\langle A\operatorname{grad}u,\nu\rangle\)，其中 \(\mu_g\) 和 \(\mu_\partial\) 分别为体积密度及边界密度。记 \(\mathcal T(\partial M)=\gamma H^1(M)\) 并配商范数，\(V^\dagger\) 为 \(V\) 的有界反线性对偶。对任意 \(F\in(H_0^1(M))^\dagger\)、\(\psi\in\mathcal T(\partial M)\)，存在唯一 \(u\in H^1(M;\mathbb C)\) 满足
	\[
		\gamma u=\psi,\qquad Q(u,v)=F(v)\quad(v\in H_0^1(M)).
	\]
	而且 \(\|u\|_{H^1}\le C(\|F\|+\|\psi\|_{\mathcal T})\)。当 \(F(v)=(f,v)_2\)、\(f\in L^2(M)\) 时，\(u\) 是
	\[
		\mathcal E(w)=\tfrac12Q(w,w)-\operatorname{Re}(f,w)_2
	\]
	在仿射空间 \(\{w\in H^1:\gamma w=\psi\}\) 上的唯一极小点。
\end{theorem}
\begin{proof}
	取最小范数提升 \(w=E\psi\)，写 \(u=w+z\)，其中 \(z\in H_0^1\)。未知量满足
	\[
		Q(z,v)=F(v)-Q(w,v)\quad(v\in H_0^1).
	\]
	右侧为有界反线性泛函。由上一引理、\(A\) 的正定性和 \(c\ge0\)，有 \(Q(v,v)\ge\alpha\|v\|_{H^1}^2\)，\(\alpha>0\)。\Cref{g3:thm:lax-milgram}给出唯一的 \(z\) 及范数界。若更换提升，仿射空间不变；两解之差在 \(H_0^1\)，对它自身测试给出零，所以最终的 \(u\) 与提升无关。

	当右端来自 \(L^2\) 时，对任意 \(h\in H_0^1\) 展开能量差，得到
	\[
		\mathcal E(u+h)-\mathcal E(u)
		=\operatorname{Re}\bigl(Q(u,h)-(f,h)_2\bigr)+\tfrac12Q(h,h)
		=\tfrac12Q(h,h).
	\]
	故 \(u\) 唯一极小。反过来，对 \(th\) 求实参数导数，再对 \(\mathrm ih\) 求导，便恢复复弱方程。
\end{proof}

例如在 \([0,L]\) 上求 \(-u''=f\)、\(u(0)=a\)、\(u(L)=b\)，可先减去线性提升 \(a+(b-a)x/L\)。剩余未知量有零端点值，梯度能量强制。存在性因而不要求 \(f\) 的积分为零；外通量由解确定，不能再任意指定。

\ProseParagraphBreak

非齐次边界提升与变分极小的基本版本可对照 Brezis 第9.5节例2、命题9.22\cite[第9.5节例2，命题9.22]{Brezis2011}。在一般 \(F\in(H_0^1)^\dagger\) 的情形，\(F\) 只定义于零边界测试空间；存在性证明使用的是上面的移位方程，无需把 \(F\) 未经选择地作用在全部仿射边界类上。

\subsection{Neumann 兼容条件及常数核}
\label{g3:subsec:1401:03}
在\term[Neumann bianjie tiaojian]{Neumann 边界条件}[Neumann boundary condition]下，指定共法向通量。若 \(Pu=f\)、\(B_Au=g\)，第一 Green 公式要求
\begin{equation}\label{g3:eq:weak-neumann-form}
	Q(u,v)=(f,v)_2+G(\gamma v)\qquad(v\in H^1(M)),
\end{equation}
其中 \(G\in\mathcal T^\dagger\) 是给定通量；当 \(g\in L^2(\partial M)\) 时，\(G(\psi)=\int g\overline\psi\,\mu_{\partial}\)。这里对全部 \(H^1\) 函数测试，正是边界通量得以确定的原因。

\begin{theorem}\label{g3:thm:weak-general-neumann}
	设 \(M\) 为非空紧光滑带边界黎曼流形，维数至少一，每个分支均有边界；\(A\) 光滑、实对称且一致正定。取 \(P=-\operatorname{div}(A\operatorname{grad})\)，即 \(c=0\)。记 \(\gamma\) 为边界迹，\(\mathcal T=\gamma H^1(M)\) 配商范数，\(\mathcal T^\dagger\) 为其有界反线性对偶。给定 \(f\in L^2(M)\)、\(G\in\mathcal T^\dagger\)，弱问题 \(\int_M\langle A\operatorname{grad}u,\operatorname{grad}v\rangle\,\mu_g=(f,v)_{L^2}+G(\gamma v)\)（\(u,v\in H^1(M)\)，对所有 \(v\) 成立）可解当且仅当对每个连通分支 \(M_j\) 有
	\begin{equation}\label{g3:eq:neumann-compatibility}
		\int_{M_j}f\,\mu_g+G(\gamma\mathbf1_{M_j})=0.
	\end{equation}
	满足 \(\int_{M_j}u\,\mu_g=0\) 的解唯一，且
	\[
		\|u\|_{H^1}\le C\bigl(\|f\|_2+\|G\|_{\mathcal T^\dagger}\bigr).
	\]
	所有其他解与它相差逐分支常数。该解满足 \(Pu=f\) 及弱共法向通量 \(\gamma_Au=G\)。
\end{theorem}
\begin{proof}
	以 \(\mathbf1_{M_j}\) 为测试，左侧为零，得到必要条件。反过来，令 \(V\subset H^1(M)\) 为逐分支零均值的闭子空间。在 \(V\) 上，\cref{g3:lem:boundary-poincare}及一致正定性使 \(Q\) 强制。右端连续，所以 Lax--Milgram 给出唯一的 \(u\in V\) 满足全部 \(V\) 中的测试方程。

	任意 \(v\in H^1\) 唯一分解为 \(v_0+\sum_jv_{M_j}\mathbf1_{M_j}\)，其中 \(v_0\in V\)。兼容条件使常数部分的右端为零，左侧也为零，故方程对任意 \(v\) 成立。两解相减并取自身测试，得到梯度为零，从而差为逐分支常数。内部测试给出分布方程，\cref{g3:prop:weak-conormal-trace}随后把剩余配对识别为 \(\gamma_Au=G\)。范数界来自强制常数与迹范数的定义。
\end{proof}

若通量由 \(g\in L^2(\partial M)\) 给出，则兼容条件是
\[
	\int_{M_j}f\,\mu_g=-\int_{\partial M_j}g\,\mu_{\partial}.
\]
负号由 \(P=-\operatorname{div}(A\operatorname{grad})\) 决定。对于区间 \([0,L]\)，\(g(0)=-u'(0)\)、\(g(L)=u'(L)\)，所以 \(\int_0^Lf+g(0)+g(L)=0\)。当 \(f=1\) 而两端外通量均为零时，条件立即失败。

\ProseParagraphBreak

若增加正零阶项，例如 \(P=\Delta_g+1\)，则 \(Q(u,u)=\|\mathrm du\|_2^2+\|u\|_2^2\) 在全部 \(H^1\) 上强制，每个 \(f,G\) 都有唯一解。此时对常数测试得到 \(\int u=\int f+G(1)\)，这是确定解的平均值的等式。对应齐次 Neumann 变分问题见 Brezis 第9.5节例4、命题9.24\cite[第9.5节例4，命题9.24]{Brezis2011}。

\subsection{混合边界条件与角点奇性}
\label{g3:subsec:1401:04}
同一边界可以在不同部分施加不同条件。设 \(\Gamma_D,\Gamma_N\subset\partial M\) 是不交的相对开集，余下交界集合的边界测度为零。定义
\[
	V_D=\{v\in H^1(M):\gamma v=0\text{ 于 }\Gamma_D\}.
\]
这是闭子空间。若 \(\Gamma_D\) 在每个连通分支边界上均有正测度，则刚才的 Poincaré 证明直接适用。给定 \(w\in H^1\)，用 \(w+V_D\) 表示指定 \(\Gamma_D\) 上函数值的仿射空间。

\begin{theorem}\label{g3:thm:weak-mixed-boundary}
	设 \(M\) 为紧光滑带边界黎曼流形，维数至少一，\(A\) 光滑、实对称且一致正定，\(c\ge0\) 实光滑。令 \(\gamma\) 为边界迹，\(\nu\) 为外单位法向，\(P=-\operatorname{div}(A\operatorname{grad})+c\)，\(Q(u,v)=\int_M(\langle A\operatorname{grad}u,\operatorname{grad}v\rangle+cu\overline v)\,\mu_g\)，\(B_Au=\langle A\operatorname{grad}u,\nu\rangle\)，其中 \(\mu_g\) 和 \(\mu_\partial\) 分别为体积密度及边界密度。设 \(\Gamma_D,\Gamma_N\) 是边界上不交的相对开集，其并的补集有边界测度零，且每个连通分支的 \(\Gamma_D\) 有正测度；令 \(V_D=\{v\in H^1:\gamma v|_{\Gamma_D}=0\}\)。对 \(f\in L^2(M)\)、\(g\in L^2(\Gamma_N)\) 和 \(w\in H^1(M)\)，存在唯一 \(u\in w+V_D\)，使
	\[
		Q(u,v)=(f,v)_2+\int_{\Gamma_N}g\,\overline{\gamma v}\,\mu_{\partial}
		\qquad(v\in V_D).
	\]
	并且 \(\|u\|_{H^1}\le C(\|f\|_2+\|g\|_{L^2(\Gamma_N)}+\|w\|_{H^1})\)。它唯一极小化同一仿射空间上的能量
	\[
		\tfrac12Q(u,u)-\operatorname{Re}(f,u)_2
		-\operatorname{Re}\int_{\Gamma_N}g\,\overline{\gamma u}\,\mu_{\partial}.
	\]
\end{theorem}
\begin{proof}
	迹估计保证边界积分是 \(V_D\) 上的连续反线性泛函。写 \(u=w+z\)，由\cref{g3:lem:boundary-poincare}，\(Q\) 在 \(V_D\) 上强制；应用 Lax--Milgram 便得到唯一 \(z\) 和所列估计。能量差的展开与 Dirichlet 情形相同，线性项由弱方程消去，留下严格正的 \(Q(h,h)/2\)。内部紧支测试给出 \(Pu=f\)；若解光滑，在 \(\Gamma_N\) 内任取紧支边界测试并提升，Green 公式给出 \(B_Au=g\)。
\end{proof}

这种\term[hunhe bianjie tiaojian]{混合边界条件}[mixed boundary condition]在能量空间中十分自然，但其交界点会改变高阶正则性。先看真正的几何角点。

\begin{example}\label{g3:exm:reentrant-corner-singularity}
	取平面扇形 \(S_\omega=\{0<r<1,\ 0<\theta<\omega\}\)，其中 \(\pi<\omega<2\pi\)，令 \(\alpha=\pi/\omega\in(1/2,1)\)。选光滑径向截断 \(\chi\)，在 \(r\le1/4\) 时等于一、在 \(r\ge1/2\) 时为零，置
	\[
		u(r,\theta)=\chi(r)r^\alpha\sin(\alpha\theta).
	\]
	在角点附近，\(\sum_i\partial_i^2u=0\)，且两条边上 \(u=0\)。一阶导数平方的径向积分为常数乘以 \(\int_0^{1/4}r^{2\alpha-1}\,\mathrm dr<\infty\)，所以 \(u\in H^1\)。在 \(r<2\varepsilon\) 切去此函数，\(H^1\) 误差为 \(O(\varepsilon^\alpha)\)；切去后的函数在光滑边界片上具有零迹，可向内平移并光滑化，故 \(u\in H_0^1(S_\omega)\)。

	设 \(f=\Delta u\)，采用本书非负号。\(f\) 只支于 \(1/4\le r\le1/2\)，在各个闭边界片光滑，因而属于 \(L^2\)。内部积分及 \(H_0^1\) 稠密性说明 \(u\) 是对应 Dirichlet 弱解。但在角点附近
	\[
		|\mathrm D^2u|^2=2\alpha^2(\alpha-1)^2r^{2\alpha-4},\qquad
		\int_0^{1/4}r^{2\alpha-3}\,\mathrm dr=\infty.
	\]
	因此 \(u\notin H^2\)。零边界、光滑右端与内部光滑性共同存在，凹角仍使整体二阶估计失败。
\end{example}

\begin{example}\label{g3:exm:smooth-boundary-mixed-singularity}
	光滑边界也会出现相同机制。在一个局部等于上半平面的光滑有界区域中，以原点为条件交界点：\(\theta=0\) 的正半轴施加 Dirichlet 条件，\(\theta=\pi\) 的负半轴施加 Neumann 条件。令
	\[
		u=\chi(r)r^{1/2}\sin(\theta/2),
	\]
	其中截断支于这个平边界图内，并在图外补零。在正半轴 \(u=0\)；在负半轴 \(\partial_\theta u=0\)，故外法向导数为零。\(u\in H^1\)，而 Hessian 平方积分含 \(\int_0r^{-2}\,\mathrm dr\)，所以 \(u\notin H^2\)。右端 \(f=\Delta u\) 仍光滑且支于远离原点的环带。

	验证它满足混合弱问题时，先以有界的 \(v\in V_D\) 测试，并用径向截断 \(\eta_\varepsilon\) 消去原点：\(\eta_\varepsilon=0\) 于 \(r<\varepsilon\)、等于一于 \(r>2\varepsilon\)、\(|\mathrm D\eta_\varepsilon|\le C/\varepsilon\)。穿过小环带的额外项满足
	\[
		\left|\int v\,\mathrm du\cdot\mathrm d\eta_\varepsilon\right|
		\le C\|v\|_\infty\varepsilon^{-1}
		\int_\varepsilon^{2\varepsilon}r^{-1/2}r\,\mathrm dr
		\le C\|v\|_\infty\varepsilon^{1/2}\longrightarrow0.
	\]
	其余项由 \(H^1\) 可积性取极限。在实、虚部分别作截断，任意 \(V_D\) 函数均可在 \(H^1\) 中由有界的 \(V_D\) 函数逼近，故弱恒等式对全部测试成立。区域的几何边界在原点光滑，失去二阶正则性的原因是边界条件在那里改变。
\end{example}

\begin{figure}[htbp]
	\centering
	% !TeX root = ../Book3.tex
\begin{tikzpicture}[scale=1.05,>=stealth]
	\fill[blue!4] (-2.7,0) rectangle (2.7,1.8);
	\draw[very thick,blue!65!black] (-2.7,0)--(0,0);
	\draw[very thick,red!65!black] (0,0)--(2.7,0);
	\fill (0,0) circle (1.6pt);
	\node[below=4pt] at (0,0) {条件交界点};
	\node[below=4pt] at (-1.8,0) {\(\Gamma_N\)};
	\node[below=4pt] at (1.8,0) {\(\Gamma_D\)};
	\draw[dashed] (1.1,0) arc[start angle=0,end angle=180,radius=1.1];
	\draw[->] (0,0)--(35:1.55) node[right] {\(r\)};
	\draw[->] (.5,0) arc[start angle=0,end angle=35,radius=.5];
	\node at (.72,.19) {\(\theta\)};
	\node at (0,1.53) {\(u=r^{1/2}\sin(\theta/2)\)};
\end{tikzpicture}

	\caption{光滑边界上的混合条件交界}
	\label{g3:fig:mixed-boundary-interface}
\end{figure}

因此下一节的整体光滑边界估计分别处理完整的 Dirichlet 条件或完整的 Neumann 条件。两类条件若只分布在互不相交的整个边界分支上，各个边界图仍只有一种条件，局部证明可以逐图拼接；一旦出现交界，就需要另行分析其奇异指数。
% 纲要条目排为编号小节标题，并进入目录与书签；说明段仍保留于纲要注释中。

% !TeX root = ../../Book3.tex
% 纲要标题：### 14.2 边界正则性
% 纲要位置：几何学计划纲要.md，第 2680 行。
% 正文状态：初稿；尚未正式审核。

% 纲要内容开始
% 
% * 边界展平与切向差商
% * 法向导数的方程恢复
% * 齐次及非齐次边界数据
% * 高阶边界估计
% * 整体估计中的边界条件
% 
% 纲要内容结束

\section{边界正则性}
\label{g3:sec:1402}

前一节的能量方法给出 \(H^1\) 弱解。要进一步求二阶导数，内部差商论证只有在平移仍留在区域中时才能使用；贴近边界时，允许的平移方向和测试函数都受边界条件约束。本节先沿边界求导，再利用椭圆方程求出尚未控制的法向导数。这一方法同时解释了边界数据为什么会出现在整体估计中。

以下以有界光滑域 \(\Omega\subset\mathbb R^n\) 上的标量算子为主：
\[
	Pu=-\partial_i(a^{ij}\partial_j u)+cu,\qquad
	B_\nu u=\nu_i a^{ij}\partial_j u.
\]
重复指标求和，\(A=(a^{ij})\) 为实对称一致正定矩阵，\(\nu\) 为外单位法向量。函数可取复值，积分配对对第一变量线性。记 \(\gamma=\operatorname{Tr}\)。正则性估计保留 \(\|u\|_{L^2}\) 项，因而不要求 \(c\ge0\)；正性在前一节负责存在性与唯一性。

\subsection{边界展平与切向差商}
\label{g3:subsec:1402:01}
取边界展平图 \(x=\Phi(y)\)，把区域映成
\[
	Q_R^+=(-R,R)^{n-1}\times(0,R),\qquad
	\Gamma_R=(-R,R)^{n-1}\times\{0\}.
\]
令 \(v=u\circ\Phi\)、\(J=\mathrm D\Phi\)、\(\rho=|\det J|\)。直接在弱形式中换元，主部系数变为
\begin{equation}\label{g3:eq:boundary-flattened-coefficients}
	\widetilde A=\rho J^{-1}(A\circ\Phi)J^{-T},\qquad
	\widetilde c=\rho(c\circ\Phi),\qquad
	\widetilde f=\rho(f\circ\Phi).
\end{equation}
图与逆图的一阶导数有界，\(\rho\) 有正下界，所以 \(\widetilde A\) 仍一致正定。\(C^2\) 边界图和 \(W^{1,\infty}\) 主系数保证 \(\widetilde A\in W^{1,\infty}\)。这里变换的是弱形式：不需要先假定 \(u\) 已有二阶导数。齐次 Dirichlet 条件成为 \(\gamma v=0\)；齐次 Neumann 条件使贴近平面 \(\Gamma_R\) 的测试函数仍可自由取值。

下面保留一个散度形式的源项。高阶求导会产生这样的项，它使后面的递推能够使用同一个局部结论。

\begin{lemma}\label{g3:lem:boundary-tangential-differences}
	设 \(n\ge2\)、\(0<r<R\)，记 \(Q_a^+=(-a,a)^{n-1}\times(0,a)\)、\(\Gamma_a=(-a,a)^{n-1}\times\{0\}\)，\(\gamma\) 为平面边界迹。设 \(A\in W^{1,\infty}(Q_R^+)\) 为实对称矩阵，且
	\(a^{ij}(y)\xi_j\xi_i\ge\lambda|\xi|^2\)，其中 \(\lambda>0\)。设 \(v\in H^1(Q_R^+)\)、\(F\in L^2(Q_R^+)\)、\(G\in H^1(Q_R^+;\mathbb C^n)\)，并有
	\begin{equation}\label{g3:eq:flat-boundary-weak-div-source}
	\int a^{ij}\partial_jv\,\overline{\partial_i\varphi}
	=\int F\overline\varphi+\int G^i\overline{\partial_i\varphi}.
	\end{equation}
	积分均在 \(Q_R^+\) 上进行；测试函数在侧面与顶面附近消失。允许下列两类边界条件之一：
	\begin{enumerate}
		\item \(\gamma v=0\) 于 \(\Gamma_R\)，且测试函数满足 \(\gamma\varphi=0\)；
		\item 对平面 \(\Gamma_R\) 上的迹不加限制，式中允许全部上述测试函数。
	\end{enumerate}
	则对每个 \(\alpha<n\)，\(\partial_\alpha\mathrm Dv\in L^2(Q_r^+)\)，且
	\[
	\|\partial_\alpha\mathrm Dv\|_{L^2(Q_r^+)}
	\le C\bigl(\|F\|_{L^2(Q_R^+)}
							+\|G\|_{H^1(Q_R^+)}+\|v\|_{H^1(Q_R^+)}\bigr).
	\]
	常数只依赖于 \(r,R,n,\lambda\) 与 \(\|A\|_{W^{1,\infty}}\)。
\end{lemma}
\begin{proof}
	记切向差商为
	\[
	\delta_h v(y)=\frac{v(y+he_\alpha)-v(y)}h.
	\]
	取实截断函数 \(\eta\)，使其在 \(Q_r^+\) 上为一，在侧面与顶面附近为零；它可在 \(\Gamma_R\) 上非零。选 \(|h|\) 小于截断支集到人工边界距离的一半，令
	\[
	z=\delta_hv,\qquad \psi=\eta^2z,\qquad
	\varphi=-\delta_{-h}\psi.
	\]
	固定 \(h\) 时这些函数属于 \(H^1\)。第一类边界条件下，切向平移保持零迹，故 \(\varphi\) 仍为许可的测试函数。迹与平移交换可先对光滑函数核对，再用 \(H^1\) 逼近和迹连续性取极限。第二类条件没有平面上的迹限制，同一个测试函数也可使用。

	差商的分部积分公式
	\(\int p\,\overline{-\delta_{-h}q}=\int\delta_hp\,\overline q\)
	与乘法公式给出
	\begin{align}
	\int a^{ij}(y+he_\alpha)\partial_jz\,
																	\overline{\partial_i(\eta^2z)}
	&=\int F\,\overline{-\delta_{-h}(\eta^2z)}
			+\int\delta_hG^i\,\overline{\partial_i(\eta^2z)}
	\notag\\
	&\quad-\int(\delta_ha^{ij})\partial_jv\,
																				\overline{\partial_i(\eta^2z)}.
	\label{g3:eq:boundary-tangential-energy}
	\end{align}
	支集留在较大的半柱内，所以没有侧面或顶面的边界项。取实部，左侧的主项至少为
	\(\lambda\|\eta\mathrm Dz\|_2^2\)；截断导数引入的项至多为
	\(C\|\eta\mathrm Dz\|_2\|z\|_2\)。

	为估计右侧，关键是让差商留在测试函数上，而不对 \(F\) 求导。由\cref{g3:lem:difference-quotient-criterion}中的差商估计，
	\[
	\|\delta_{-h}(\eta^2z)\|_2
	\le\|\partial_\alpha(\eta^2z)\|_2
	\le C\bigl(\|\eta\mathrm Dz\|_2+\|z\|_2\bigr).
	\]
	其余两项分别使用
	\[
	\|\delta_hG\|_2\le\|\partial_\alpha G\|_2,\qquad
	\|\delta_hA\|_\infty\le\|\partial_\alpha A\|_\infty,\qquad
	\|z\|_2\le C\|\mathrm Dv\|_{L^2(Q_R^+)}.
	\]
	这些范数取在包含全部平移支集的半柱上。于是记
	\(K=\|F\|_2+\|G\|_{H^1}+\|v\|_{H^1}\)，便有
	\[
	\lambda\|\eta\mathrm Dz\|_2^2
	\le CK\|\eta\mathrm Dz\|_2+CK^2
	\le \frac{\lambda}{2}\|\eta\mathrm Dz\|_2^2+C'K^2.
	\]
	吸收后得到与 \(h\) 无关的界。

	因此 \(\delta_h(\partial_jv)\) 在 \(L^2(Q_r^+)\) 中有弱收敛子列。对内部紧支光滑函数 \(\chi\) 使用差商分部积分并令 \(h\to0\)，极限满足
	\[
	\int w\,\overline\chi=-\int\partial_jv\,\overline{\partial_\alpha\chi},
	\]
	故 \(w=\partial_\alpha\partial_jv\)。弱下半连续性给出所列估计。虽然辨认弱导数只使用内部测试函数，所得积分界覆盖整个 \(Q_r^+\)，包括任意接近平面的部分。
\end{proof}

\subsection{法向导数的方程恢复}
\label{g3:subsec:1402:02}
切向估计已经控制所有含至少一个切向指标的二阶导数。唯一遗漏的是 \(\partial_n^2v\)。它不能靠向外平移取得，但方程在法向上具有非零的最高阶系数。

\begin{lemma}\label{g3:lem:flat-boundary-h2-div-source}
	设 \(n\ge2\)、\(0<r<R\)，\(Q_a^+=(-a,a)^{n-1}\times(0,a)\)。设 \(A\in W^{1,\infty}(Q_R^+;\mathbb R^{n\times n})\) 实对称且一致正定，\(F\in L^2(Q_R^+)\)、\(G\in H^1(Q_R^+;\mathbb C^n)\)、\(v\in H^1(Q_R^+)\)。假定 \(\int A\mathrm Dv\cdot\overline{\mathrm D\varphi}=\int F\overline\varphi+G\cdot\overline{\mathrm D\varphi}\) 对侧面及顶面附近消失的测试函数成立，并且底面上或者 \(v\) 及测试函数都为零迹，或者允许所有测试函数的底面迹。则
	\(v\in H^2(Q_r^+)\)，并且
	\[
	\|v\|_{H^2(Q_r^+)}
	\le C\bigl(\|F\|_{L^2(Q_R^+)}
							+\|G\|_{H^1(Q_R^+)}+\|v\|_{H^1(Q_R^+)}\bigr).
	\]
	一维半区间上相同结论成立，且不需要切向步骤。
\end{lemma}
\begin{proof}
	先把切向估计应用于一个严格介于 \(Q_r^+\) 与 \(Q_R^+\) 之间的半柱。在分布意义下，弱方程给出
	\[
	\partial_n(a^{nn}\partial_nv)
	=-F+\partial_iG^i
							-\sum_{(i,j)\ne(n,n)}\partial_i(a^{ij}\partial_jv).
	\]
	右侧属于 \(L^2\)：最后一个和式中的每个二阶导数都含切向指标，系数导数与 \(\mathrm Dv\) 的乘积也属于 \(L^2\)。因此 \(a^{nn}\partial_nv\) 有 \(L^2\) 的法向弱导数。由于
	\(a^{nn}\ge\lambda\)，其倒数属于 \(W^{1,\infty}\)。弱乘法公式遂给出
	\[
	\partial_n^2v
	=\frac1{a^{nn}}\partial_n(a^{nn}\partial_nv)
			-\frac{\partial_na^{nn}}{a^{nn}}\partial_nv\in L^2.
	\]
	逐项估计即得结论。这个写法先对通量 \(a^{nn}\partial_nv\) 求导，再除以系数，因而不预先假定 \(\partial_n^2v\) 是函数。一维时和式为空，同一通量计算直接完成证明。
\end{proof}

\begin{theorem}\label{g3:thm:boundary-h2-homogeneous}
	设 \(\Omega\subset\mathbb R^n\) 为有界 \(C^2\) 域，\(A\in W^{1,\infty}(\Omega)\) 为实对称一致正定矩阵，\(c\in L^\infty(\Omega;\mathbb R)\)，
	\(P=-\partial_i(a^{ij}\partial_j)+c\)。记 \(\gamma\) 为边界迹，\(B_\nu u=\nu_i a^{ij}\partial_ju\) 为外共法向通量。若 \(f\in L^2(\Omega)\)，而 \(u\in H^1(\Omega)\) 满足下列齐次边值问题之一：
	\[
	Pu=f,\quad\gamma u=0;\qquad
	Pu=f,\quad B_\nu u=0,
	\]
	其中第一式以全部 \(H_0^1(\Omega)\) 函数测试，第二式以全部 \(H^1(\Omega)\) 函数测试，则
	\(u\in H^2(\Omega)\)，且
	\begin{equation}\label{g3:eq:homogeneous-boundary-h2-estimate}
	\|u\|_{H^2(\Omega)}
	\le C\bigl(\|f\|_{L^2(\Omega)}+\|u\|_{L^2(\Omega)}\bigr).
	\end{equation}
	Neumann 条件在升级后等于 \(H^2\) 函数的通常共法向导数迹。
\end{theorem}
\begin{proof}
	在边界图中按式\eqref{g3:eq:boundary-flattened-coefficients}换元，将零阶项移到右侧，取
	\(F=\widetilde f-\widetilde c\,v\)、\(G=0\)。前两个引理给出较小边界图上的 \(H^2\) 估计。内部图使用\cref{g3:thm:weak-interior-h2}；其假设由实对称一致正定主部满足。有限覆盖和坐标范数等价得到
	\[
	\|u\|_{H^2}\le C(\|f\|_2+\|u\|_{H^1}).
	\]
	两类齐次问题都允许以 \(u\) 本身测试。取实部并记 \(c_-=\max\{-c,0\}\)，得
	\[
	\lambda\|\mathrm Du\|_2^2
	\le \|f\|_2\|u\|_2+\|c_-\|_\infty\|u\|_2^2
	\le C(\|f\|_2^2+\|u\|_2^2).
	\]
	代回即得整体估计。

	升级为 \(H^2\) 后，\(A\mathrm Du\in H^1\)，故其边界迹存在。由\cref{g3:prop:weak-boundary-integration}逐分量分部积分，
	\[
	\int_\Omega\bigl(a^{ij}\partial_ju\,\overline{\partial_iw}
																									+cu\overline w\bigr)
	-\int_\Omega f\overline w
	=\int_{\partial\Omega}B_\nu u\,\overline{\gamma w}
	\quad(w\in H^1).
	\]
	左侧在齐次 Neumann 情形为零。每个光滑边界函数都有光滑的局部法向常值延拓，拼接后可作为 \(\gamma w\)，所以 \(B_\nu u=0\) 的迹成立。这也确认了它与前一节定义的弱共法向导数一致。
\end{proof}

在紧的光滑带边界 黎曼流形上，上述定理同样适用于
\(P=-\operatorname{div}(A\nabla)+c\)，其中 \(A\) 是切丛上光滑、实对称、一致正定的自同态。坐标中的体积密度并入 \(\widetilde A,\widetilde c,\widetilde f\)，正密度与有限坐标覆盖保持所有估计；边界积分使用诱导的面积密度。因此以后在流形上应用本节的标量结论时，共法向导数仍取 \(A\nabla u\) 与几何外法向 \(\nu\) 的配对。

\subsection{齐次及非齐次边界数据}
\label{g3:subsec:1402:03}
若边界值不为零，可以减去一个具有这些边界值的函数。但为了得到 \(H^2\) 结论，减去的函数也必须属于 \(H^2\)。这一要求把 Dirichlet 数据放在 \(H^{3/2}(\partial\Omega)\)，把 Neumann 数据放在 \(H^{1/2}(\partial\Omega)\)。先建立这些空间及提升，才能准确实施这一步。

设 \(m\ge1\)。对 \(s\ge0\)，在 \(\mathbb R^m\) 上用单位化 Fourier 变换定义
\[
	\|h\|_{H^s(\mathbb R^m)}^2
	=\int_{\mathbb R^m}\langle\xi\rangle^{2s}|\widehat h(\xi)|^2\,\mathrm d\xi,
	\qquad \langle\xi\rangle=(1+|\xi|^2)^{1/2}.
\]
整数阶时它与弱导数范数等价。本节只需要半整数阶。为把定义搬到边界，须先确认它在换坐标及乘截断下稳定。

\begin{lemma}\label{g3:lem:half-integer-localization}
	设 \(m\ge1\)，对整数 \(\ell\ge0\)，\(H^{\ell+1/2}(\mathbb R^m)\) 的范数等价于
	\[
	\|h\|_{H^\ell}^2+
	\sum_{|\alpha|=\ell}
	\iint_{\mathbb R^m\times\mathbb R^m}
	\frac{|\partial^\alpha h(x)-\partial^\alpha h(y)|^2}
						{|x-y|^{m+1}}\,\mathrm dx\,\mathrm dy.
	\]
	光滑紧支函数的乘法，以及紧支局部函数在光滑坐标变换下的拉回，均在这些空间上有界。因此，对任意有界光滑域 \(\Omega\subset\mathbb R^{m+1}\)，可用有限边界图和单位分解定义
	\(H^{\ell+1/2}(\partial\Omega)\)，且不同选择给出等价范数。
\end{lemma}
\begin{proof}
	对 Schwartz 函数 \(h\)，将 \(y=x+z\) 代入双重积分，使用 Plancherel 等式与非负积分换序，得到
	\[
	[h]_{1/2}^2
	=\int|\widehat h(\xi)|^2
			\left(\int\frac{|e^{i\xi\cdot z}-1|^2}{|z|^{m+1}}\,\mathrm dz\right)\mathrm d\xi
	=C_m\int|\xi|\,|\widehat h(\xi)|^2\,\mathrm d\xi.
	\]
	内层积分在零点附近由分子 \(O(|z|^2)\) 控制，在无穷远由分子有界控制；旋转和伸缩给出最后的等式，且 \(0<C_m<\infty\)。于是
	\(\|h\|_2^2+[h]_{1/2}^2\) 等价于 \(H^{1/2}\) 范数。对 \(\ell\) 阶导数求和，并使用
	\(\sum_{|\alpha|=\ell}|\xi^\alpha|^2\asymp|\xi|^{2\ell}\)，得到一般等价式。在 Fourier 侧用光滑紧支函数于加权 \(L^2\) 中逼近，其逆变换为 Schwartz 函数。对逼近列应用上述等式并完成，便得 Fourier 定义与双重积分范数的等价性。

	若 \(a\) 是有界 Lipschitz 函数，则
	\[
	|a(x)h(x)-a(y)h(y)|
	\le \|a\|_\infty|h(x)-h(y)|
						+|a(x)-a(y)|\,|h(y)|.
	\]
	第二项在 \(|x-y|<1\) 用 Lipschitz 界，在 \(|x-y|\ge1\) 用有界性；相应核积分均有限。故乘以 \(a\) 在 \(H^{1/2}\) 上有界。高阶情形对乘积求 \(\ell\) 阶导数：每一项是光滑系数乘某个 \(\partial^\beta h\)，\(|\beta|\le\ell\)，而 Fourier 权重给出
	\(\|\partial^\beta h\|_{H^{1/2}}\le C\|h\|_{H^{\ell+1/2}}\)。

	对坐标变换，在包含局部支集的较大坐标邻域上，映射与逆映射均为 Lipschitz，体积比有上下界，因而双重积分的距离核和体积因子受常数控制。两个点中有一个离开该较大邻域时，紧支集与其补集之间有正距离，所得积分只需 \(L^2\) 控制。高阶拉回的链式法则把问题归结为已证的 \(H^{1/2}\) 拉回与光滑乘法。两组边界图取交并细分单位分解，便得两个全局范数的相互控制。局部卷积再拼接还说明 \(C^\infty(\partial\Omega)\) 在这些空间中稠密。
\end{proof}

\begin{theorem}\label{g3:thm:half-integer-boundary-trace}
	设 \(\Omega\subset\mathbb R^n\) 为有界光滑域，\(q\ge1\) 为整数。迹映射有连续满射延拓
	\[
	\gamma:H^q(\Omega)\longrightarrow H^{q-1/2}(\partial\Omega),
	\]
	并有有界线性右逆。若 \(q\ge2\)，联合迹
	\[
	(\gamma,\gamma\partial_\nu):H^q(\Omega)\longrightarrow
	H^{q-1/2}(\partial\Omega)\times H^{q-3/2}(\partial\Omega)
	\]
	也有有界线性右逆。这里 \(\partial_\nu\) 表示普通外法向导数。这些结论对紧的光滑带边界流形及光滑向量丛逐分量成立；丛情形固定光滑联络，并将法向导数换为 \(\nabla_\nu\)。
\end{theorem}
\begin{proof}
	先在 \(\mathbb R^{n-1}\times(0,\infty)\) 上沿 \(y'\) 作部分 Fourier 变换，记 \(t=y_n\)、\(\lambda_\xi=\langle\xi\rangle\)。对闭半空间上的光滑紧支函数，一维分部积分给出
	\[
	|\widehat v(\xi,0)|^2
	\le2\int_0^\infty|\widehat v|\,|\partial_t\widehat v|\,\mathrm dt
	\le\lambda_\xi\int_0^\infty|\widehat v|^2\,\mathrm dt
					+\lambda_\xi^{-1}\int_0^\infty|\partial_t\widehat v|^2\,\mathrm dt.
	\]
	两边乘 \(\lambda_\xi^{2q-1}\) 并对 \(\xi\) 积分，右侧由整数阶 \(H^q\) 范数控制。这证明 \(\gamma:H^q\to H^{q-1/2}\) 的界。将 \(q\) 换为 \(q-1\) 并应用于 \(\partial_tv\)，便得一阶法向迹的界。闭半空间上的高阶光滑稠密性已在\cref{g3:prop:higher-order-boundary-extension}中用向内平移与卷积证明，故两种迹均连续延拓。

	右逆可直接写出。只指定零阶迹 \(g\) 时，取
	\[
	\widehat{E_0g}(\xi,t)=e^{-t\lambda_\xi}\widehat g(\xi).
	\]
	同时指定 \(v|_{t=0}=g\)、\(\partial_tv|_{t=0}=h\) 时，取
	\begin{equation}\label{g3:eq:fourier-boundary-jet-lift}
	\widehat{E(g,h)}(\xi,t)
	=e^{-t\lambda_\xi}
							\bigl((1+t\lambda_\xi)\widehat g(\xi)+t\widehat h(\xi)\bigr).
	\end{equation}
	因为 \(\partial_t[(1+t\lambda_\xi)e^{-t\lambda_\xi}]_{t=0}=0\)，两个迹恰为 \(g,h\)。对任意 \(|\alpha|+b\le q\)，求 \(\partial_{y'}^\alpha\partial_t^b\) 后，\(g\) 项由常数倍
	\(\lambda_\xi^{|\alpha|+b}(1+t\lambda_\xi)e^{-t\lambda_\xi}|\widehat g|\)
	控制，\(h\) 项由常数倍
	\(\lambda_\xi^{|\alpha|+b-1}(1+t\lambda_\xi)e^{-t\lambda_\xi}|\widehat h|\)
	控制。变量代换 \(s=t\lambda_\xi\) 因而给出
	\[
	\|E(g,h)\|_{H^q}^2
	\le C_q\left(
			\int\lambda_\xi^{2q-1}|\widehat g|^2\,\mathrm d\xi+
			\int\lambda_\xi^{2q-3}|\widehat h|^2\,\mathrm d\xi\right).
	\]
	\(E_0\) 有相同的第一项估计。先对光滑数据计算，再按此界完成，便得到全部数据的线性提升；迹连续性保证完成后的迹仍正确。

	现在用边界法向领圈拼接。欧氏情形的领圈映射为
	\(\Phi(p,t)=p-t\nu(p)\)。在 \(t=0\) 它的微分可逆；局部逆函数定理和边界紧性给出统一的正宽度。若任意小宽度仍有两个不同点映到同一点，取收敛子列后它们趋向同一边界点，便与该点附近的局部单射性矛盾。因此可取一个全局领圈。

	取边界单位分解 \(\{\psi_a\}\)，把 \(\psi_ag,\psi_ah\) 在边界图中视为紧支函数。在领圈坐标里 \(\partial_t=-\partial_\nu\)，故对要求的外法向数据使用式\eqref{g3:eq:fourier-boundary-jet-lift}中的第二输入 \(-\psi_ah\)。将提升乘以切向截断 \(\theta_a\) 及法向截断 \(\eta(t)\)：\(\theta_a\) 在 \(\operatorname{supp}\psi_a\) 附近为一，\(\eta\) 在 \(t=0\) 附近恒为一。送回领圈并在其外补零，有限求和便得全局提升。截断不改变两个指定迹，也不损失整数阶范数；\cref{g3:lem:half-integer-localization}控制输入的局部范数。只需零阶迹时用 \(E_0\) 作同样拼接。

	一般流形用由法向测地线给出的领圈，局部可逆性和紧性论证相同。向量丛在领圈上取沿法向平行移动的局部标架，逐分量构造再用单位分解求和；变换矩阵的光滑乘法有界性给出范数界。一维时边界是有限个端点，所有列出的边界空间都取端点数据的有限维范数；以局部常数、一次函数乘截断即可分别指定函数值与外法向导数。
\end{proof}

特别地，前一节定义的迹像空间有具体描述
\[
	\mathcal T(\partial\Omega)=\gamma H^1(\Omega)=H^{1/2}(\partial\Omega),
	\qquad
	\inf_{\gamma w=g}\|w\|_{H^1}\asymp\|g\|_{H^{1/2}}.
\]
一个方向由迹估计得到，另一个方向由有界右逆得到。因此前一节的连续反线性对偶 \(\mathcal T^\dagger\) 也可记作 \(H^{-1/2}(\partial\Omega)\)。这个负阶记号表示对偶空间，不要求对边界分布另作 Fourier 定义。
\symidx{\(H^{q-1/2}(\partial\Omega)\)：以局部 Fourier 范数定义的半整数阶边界 Sobolev 空间}

\begin{corollary}\label{g3:cor:boundary-conormal-lifting}
	设 \(\Omega\subset\mathbb R^n\) 为有界光滑域，\(A\) 光滑、实对称且一致正定；记 \(\gamma\) 为边界迹，\(B_\nu u=\nu_i a^{ij}\partial_ju\)，\(\nu\) 为外单位法向。对每个整数 \(q\ge2\)，存在有界线性算子
	\[
	E_N:H^{q-3/2}(\partial\Omega)\longrightarrow H^q(\Omega),
	\qquad \gamma E_Ng=0,\quad B_\nu E_Ng=g.
	\]
\end{corollary}
\begin{proof}
	令 \(a_\nu=\nu\cdot A\nu\)，它是光滑正函数，有正下界。对联合迹使用前一定理，指定
	\(\gamma w=0\)、\(\gamma\partial_\nu w=g/a_\nu\)。光滑乘法在半整数空间上有界，故提升有界。又因零阶迹为零，所有切向一阶导数的迹为零，从而
	\(B_\nu w=a_\nu\,\gamma\partial_\nu w=g\)。
	这里迹与切向求导的交换先在光滑函数上成立，再由 \(H^q\) 光滑逼近及两个迹映射的连续性成立。
\end{proof}

\begin{corollary}\label{g3:cor:boundary-h2-inhomogeneous}
	设 \(\Omega\) 为有界光滑域，\(A,c\) 光滑，\(A\) 实对称一致正定，\(c\) 实值。令 \(P=-\partial_i(a^{ij}\partial_j)+c\)，\(\gamma\) 为迹，\(B_\nu u=\nu_i a^{ij}\partial_ju\)；\(Q(u,w)=\int_\Omega(a^{ij}\partial_ju\overline{\partial_iw}+cu\overline w)\,\mathrm dx\)。若 \(f\in L^2(\Omega)\)，则满足 \(Pu=f\) 的 \(H^1\) 弱解在下列任一条件下属于 \(H^2(\Omega)\)：
	\[
	\gamma u=g_D\in H^{3/2}(\partial\Omega),\qquad
	B_\nu u=g_N\in H^{1/2}(\partial\Omega).
	\]
	分别有
	\[
	\begin{split}
	\|u\|_{H^2}&\le C\bigl(\|f\|_2+\|g_D\|_{H^{3/2}}+\|u\|_2\bigr),\\
	\|u\|_{H^2}&\le C\bigl(\|f\|_2+\|g_N\|_{H^{1/2}}+\|u\|_2\bigr).
	\end{split}
	\]
	Neumann 弱条件是
	\(Q(u,w)=(f,w)_{L^2}+\int_{\partial\Omega}g_N\overline{\gamma w}\)
	对全部 \(w\in H^1\) 成立。
\end{corollary}
\begin{proof}
	Dirichlet 情形取有界 \(H^2\) 提升 \(w=E_0g_D\)；Neumann 情形取 \(w=E_Ng_N\)。则 \(u-w\) 满足相应的齐次弱条件和方程
	\(P(u-w)=f-Pw\)。对它应用\cref{g3:thm:boundary-h2-homogeneous}，再用
	\(\|Pw\|_2\le C\|w\|_{H^2}\) 及提升界，即得结论。
\end{proof}

\subsection{高阶边界估计}
\label{g3:subsec:1402:04}
继续求导时，法向导数仍由方程恢复，边界条件仍只由切向求导来保持。为了不在 Neumann 情形遗漏系数导数所产生的边界项，应始终对弱形式求切向导数；这样交换子自然进入式\eqref{g3:eq:flat-boundary-weak-div-source}的向量源项。

\begin{theorem}\label{g3:thm:boundary-higher-regularity}
	设 \(\Omega\subset\mathbb R^n\) 为有界光滑域，\(A\in C^\infty(\overline\Omega)\) 为实对称一致正定矩阵，\(c\in C^\infty(\overline\Omega;\mathbb R)\)，
	\(P=-\partial_i(a^{ij}\partial_j)+c\)。记 \(\gamma\) 为边界迹，\(B_\nu u=\nu_i a^{ij}\partial_ju\)，\(\nu\) 为外单位法向。令 \(k\ge0\) 为整数，\(f\in H^k(\Omega)\)。若 \(u\in H^1(\Omega)\) 是满足下列任一条件的弱解，则 \(u\in H^{k+2}(\Omega)\)：
	\[
	Pu=f,\quad\gamma u=g_D\in H^{k+3/2}(\partial\Omega);
	\qquad
	Pu=f,\quad B_\nu u=g_N\in H^{k+1/2}(\partial\Omega).
	\]
	相应估计为
	\begin{align}
	\|u\|_{H^{k+2}}
	&\le C_k\bigl(\|f\|_{H^k}+\|g_D\|_{H^{k+3/2}}
																									+\|u\|_{L^2}\bigr),\label{g3:eq:dirichlet-higher-boundary-estimate}\\
	\|u\|_{H^{k+2}}
	&\le C_k\bigl(\|f\|_{H^k}+\|g_N\|_{H^{k+1/2}}
																									+\|u\|_{L^2}\bigr).\label{g3:eq:neumann-higher-boundary-estimate}
	\end{align}
	常数由区域的有限个边界图、椭圆性下界和有限阶系数界控制。相同结论适用于紧的光滑带边界 黎曼流形上的标量算子
	\(-\operatorname{div}(A\nabla)+c\)。
\end{theorem}
\begin{proof}
	先证齐次条件。\(k=0\) 已证。设 \(k\ge1\)，并已由前一阶结论知道
	\(u\in H^{k+1}\)。在展平坐标中仍将零阶项移到右侧，记主系数为 \(a^{ij}\)，记 \(F=\widetilde f-\widetilde c\,v\)。于是 \(F\in H^k\)，且
	\(\|F\|_{H^k}\le C(\|f\|_{H^k}+\|u\|_{H^k})\)。

	对纯切向多重指标 \(\beta\)，\(|\beta|=k\)，令 \(w=\partial_{y'}^\beta v\in H^1\)。在原弱等式中对测试函数作切向分部积分，得到
	\begin{equation}\label{g3:eq:boundary-tangential-commutator}
	\int a^{ij}\partial_jw\,\overline{\partial_i\varphi}
	=\int\partial_{y'}^\beta F\,\overline\varphi
								+\int G_\beta^i\,\overline{\partial_i\varphi},
	\qquad
	G_\beta^i=
	-\sum_{0<\sigma\le\beta}\binom\beta\sigma
			(\partial_{y'}^\sigma a^{ij})
							\partial_j\partial_{y'}^{\beta-\sigma}v.
	\end{equation}
	这一等式可先以光滑测试函数的切向导数代入原弱形式，再按分布导数的定义移动导数；因此不要求 \(v\) 预先具有 \(k+2\) 阶导数。在 Dirichlet 情形，
	\(\gamma w=\partial_{y'}^\beta\gamma v=0\)，其迹交换由前一小节的迹定理和 \(H^{k+1}\) 逼近给出；在 Neumann 情形，测试函数在平面上仍任意，交换子已包含在 \(G_\beta\) 中。

	每一项 \(G_\beta\) 对 \(v\) 至多使用 \(k\) 阶导数，对它再求一个导数至多使用 \(k+1\) 阶，故
	\[
	\|G_\beta\|_{H^1}\le C\|v\|_{H^{k+1}}.
	\]
	\Cref{g3:lem:flat-boundary-h2-div-source}应用于 \(w\)，给出
	\[
	\|\partial_{y'}^\beta v\|_{H^2(Q_r^+)}
	\le C\bigl(\|F\|_{H^k(Q_R^+)}
																									+\|v\|_{H^{k+1}(Q_R^+)}\bigr).
	\]
	因而所有总阶数为 \(k+2\)、法向指标个数不超过二的导数都已受控。

	还需恢复法向指标更多的导数。写成非散度形式，
	\begin{equation}\label{g3:eq:normal-high-derivative-recovery}
	-a^{nn}\partial_n^2v
	=F+\sum_{(i,j)\ne(n,n)}a^{ij}\partial_i\partial_jv
							+\sum_{i,j}(\partial_i a^{ij})\partial_jv.
	\end{equation}
	对法向指标数 \(r=3,\ldots,k+2\) 逐次归纳。给定 \(|\gamma|=k+2\)、\(\gamma_n=r\)，对上式施加 \(\partial^{\gamma-2e_n}\)，把左侧最高阶项
	\(-a^{nn}\partial^\gamma v\) 单独留下。右侧中落在 \(v\) 上的 \(k+2\) 阶导数，来自 \((i,j)\ne(n,n)\)，所以法向指标数至多为 \(r-1\)，已由归纳受控；凡有导数落在系数上的项，对 \(v\) 至多求 \(k+1\) 阶导数。对 \(F\) 只求 \(k\) 阶导数。故右侧均属于 \(L^2\)。除以 \(a^{nn}\ge\lambda\)，便得到 \(\partial^\gamma v\in L^2\) 及同型估计。所有等式在分布意义下进行，系数光滑保证求导与乘法合法。

	在内部图中，相同的弱形式求导可使用全部方向；对每个 \(|\beta|=k\) 应用内部差商证明，便直接得到 \(H^{k+2}\) 控制。该证明对 \(G_\beta\in H^1\) 的项使用
	\(\|\delta_hG_\beta\|_2\le\|\mathrm DG_\beta\|_2\)，与式\eqref{g3:eq:boundary-tangential-energy}完全相同，且无平面上的测试限制。有限个缩小的图仍覆盖闭包，于是
	\[
	\|u\|_{H^{k+2}}\le C_k(\|f\|_{H^k}+\|u\|_{H^{k+1}}).
	\]
	前一阶的齐次估计把 \(\|u\|_{H^{k+1}}\) 控制为
	\(C(\|f\|_{H^{k-1}}+\|u\|_2)\)，完成齐次归纳。一维时直接对通量方程递推，或只用式\eqref{g3:eq:normal-high-derivative-recovery}，即可得到相同结论。

	对非齐次数据，在\cref{g3:thm:half-integer-boundary-trace,g3:cor:boundary-conormal-lifting}中取 \(q=k+2\)。相应提升 \(w\) 属于 \(H^{k+2}\)，且
	\[
	\|Pw\|_{H^k}\le C_k\|w\|_{H^{k+2}}.
	\]
	将齐次结论用于 \(u-w\)，再代入提升界，分别得到两种估计。流形情形的密度、坐标变换及有限求和与二阶证明相同。
\end{proof}

若区域、系数、\(f\) 和相应边界数据都光滑，则弱解属于每个整数阶 \(H^q(\Omega)\)。\Cref{g3:prop:higher-order-boundary-extension}给出各阶有界延拓；再用高阶 Sobolev 嵌入，得到 \(u\in C^\infty(\overline\Omega)\)。这里的光滑性包含所有阶导数连续延伸到边界，正是局部内部正则性尚未提供的结论。

\subsection{整体估计中的边界条件}
\label{g3:subsec:1402:05}
整体估计包含两个不同的信息来源：方程控制内部导数，边界条件控制贴近边界的行为。下面的调和函数族可以把二者分开观察。

\begin{example}\label{g3:exm:harmonic-boundary-growth}
	在单位圆盘 \(D\subset\mathbb R^2\) 上令
	\(u_j(r,\theta)=r^j\cos(j\theta)\)，\(j\ge2\)。它们满足
	\(\Delta u_j=0\)，但
	\[
	\|u_j\|_{L^2(D)}^2=\frac{\pi}{2j+2},\qquad
	\|\nabla u_j\|_{L^2(D)}^2=\pi j,\qquad
	\sum_{a,b=1}^2\|\partial_a\partial_bu_j\|_{L^2(D)}^2
							=2\pi j^2(j-1).
	\]
	前两式由极坐标积分得到。对第三式，写 \(u_j=\operatorname{Re}(x+iy)^j\)，则
	\[
	u_{j,xx}=j(j-1)\operatorname{Re}(x+iy)^{j-2},\quad
	u_{j,xy}=-j(j-1)\operatorname{Im}(x+iy)^{j-2},\quad
	u_{j,yy}=-u_{j,xx},
	\]
	平方求和再积分即得。因此不存在对所有光滑 \(u\) 都成立的估计
	\[
	\|u\|_{H^2(D)}\le C(\|\Delta u\|_{L^2(D)}+\|u\|_{L^2(D)}).
	\]
	这些函数的 Dirichlet 数据为 \(\cos(j\theta)\)，外法向 Neumann 数据为 \(j\cos(j\theta)\)，并不满足齐次条件。高频增长来自边界数据，而非内部源项。
\end{example}

对固定的 \(0<\rho<1\)，同一族在较小圆盘 \(D_\rho\) 上的各阶导数带有 \(\rho^{j-s}\) 因子，因而与内部估计完全相容。差别在于内部图与边界之间始终有正距离。用越来越贴近边界的内部图穷尽区域时，截断函数的导数及估计常数会增长，不能由此得到一个统一的整体常数。

\ProseParagraphBreak
本节的高阶结论采用光滑边界与光滑系数。前一节若在同一连通边界上让 Dirichlet 条件与 Neumann 条件相接，交界处的测试空间不再具有本节的纯 Dirichlet 或纯 Neumann 形式；此处的论证只能应用于远离交界的边界片。下一节的微分形式边界条件则是在每个边界点分拆不同分量，所需系统估计将另行说明其测试空间和低阶边界项。
% 纲要条目排为编号小节标题，并进入目录与书签；说明段仍保留于纲要注释中。

% !TeX root = ../../Book3.tex
% 纲要标题：### 14.3 系统与补条件
% 纲要位置：几何学计划纲要.md，第 2688 行。
% 正文状态：初稿；尚未正式审核。

% 纲要内容开始
% 
% * 半空间常系数边界模型
% * 椭圆补条件的含义
% * Hodge 的绝对与相对边界条件
% * Laplace 型算子的混合边值问题
% 
% 纲要内容结束

\section{系统与补条件}
\label{g3:sec:1403}

对一个丛值方程，边界上应规定哪些分量的值，哪些分量的法向导数？仅按分量数配齐条件还不够：边界条件可能恰好消去不了某些向内部衰减的高频解。本节先在半空间把这个现象化为有限维线性代数，再证明一类具有光滑纤维分裂的边值问题。微分形式的绝对与相对边界条件属于这一类；它们的能量恒等式同时说明了边界曲率怎样进入一阶估计。

\subsection{半空间常系数边界模型}
\label{g3:subsec:1403:01}
在半空间 \(\mathbb R^{n-1}\times(0,\infty)\) 中，记坐标为 \((y,t)\)，其中 \(t\) 指向内部，所以外法向为 \(-\partial_t\)。考虑取值于 \(\mathbb C^r\) 的非负 Laplace 算子
\[
	P_0=-\partial_t^2-\sum_{a=1}^{n-1}\partial_{y_a}^2.
\]
把一个切向振荡写成 \(u(y,t)=\mathrm e^{\mathrm i y\cdot\xi}v(t)\)，其中 \(\xi\ne0\)，方程 \(P_0u=0\) 化为
\[
	-v''+|\xi|^2v=0.
\]
每个解都是 \(\mathrm e^{-|\xi|t}c+\mathrm e^{|\xi|t}d\)。要求它沿法向趋零，就剩下 \(r\) 维的解空间
\begin{equation}\label{g3:eq:decaying-boundary-model}
	v(t)=\mathrm e^{-|\xi|t}c,\qquad c\in\mathbb C^r.
\end{equation}
高切向频率对应厚度约为 \(|\xi|^{-1}\) 的边界层。能否从边界数据确定这个层内的振幅 \(c\)，是边界正则性的一个局部检验。

\begin{example}\label{g3:exm:boundary-dirichlet-robin-symbol}
	取正交分解 \(\mathbb C^r=E_{\mathrm D}\oplus E_{\mathrm N}\)，相应投影记为 \(\Pi_{\mathrm D},\Pi_{\mathrm N}\)。对 Dirichlet 与 Robin\footnote{Victor Gustave Robin（1855--1897），法国数学家，曾在 Paris 科学学院任教，研究数学分析与数学物理。} 型条件
	\[
		\Pi_{\mathrm D}v(0)=h_{\mathrm D},\qquad
		-\Pi_{\mathrm N}v'(0)+S\Pi_{\mathrm N}v(0)=h_{\mathrm N},
	\]
	其中 \(S:E_{\mathrm N}\to E_{\mathrm N}\) 是零阶项，取最高阶边界部分后，式\eqref{g3:eq:decaying-boundary-model}给出的映射是
	\[
		c\longmapsto\bigl(\Pi_{\mathrm D}c,\ |\xi|\Pi_{\mathrm N}c\bigr).
	\]
	它对每个 \(\xi\ne0\) 都可逆。纯 Dirichlet 与纯 Neumann 是两个端点；零阶矩阵 \(S\) 不改变这个最高阶检验，但会改变全局的核与特征值。
\end{example}

对标量实系数斜导数 \(b_\nu\partial_\nu+b_T\cdot\partial_y\)，相应振幅为
\[
	\bigl(b_\nu|\xi|+\mathrm i b_T\cdot\xi\bigr)c.
\]
若 \(b_\nu\ne0\)，实部已经排除非零核。若只保留切向导数，则需要重新检查：例如当 \(n\ge3\) 时，条件 \(\partial_{y_1}u=0\) 无法排除 \(\xi_1=0\) 方向的非零边界层。这里起作用的是方向，而不只是条件的条数。

\subsection{椭圆补条件的含义}
\label{g3:subsec:1403:02}
一般二阶系统也可在边界点冻结最高阶系数，沿切向作 Fourier 变换。它给出一个常系数法向常微分方程；内部椭圆性保证其特征根不落在实的法向频率轴上，但边界条件还必须控制向内部衰减的那一半解。

\begin{definition}\label{g3:def:complementing-boundary-condition}
	设 \(P\) 是作用于秩为 \(r\) 的复向量丛的光滑二阶椭圆算子。在一个边界点展平边界，并将 \(t>0\) 取为内部。对每个非零切向余向量 \(\xi\)，设冻结的主部方程
	\[
		p_2(y,0;\xi,-\mathrm i\partial_t)v=0
	\]
	的指数衰减解空间为 \(\mathcal E^+(y,\xi)\)，并假定它的维数为 \(r\)。设所取边界算子的主部在该空间上定义了总计 \(r\) 个标量条件。若其边界数据映射
	\[
		\mathcal E^+(y,\xi)\longrightarrow\mathbb C^r
	\]
	对每个边界点及每个 \(\xi\ne0\) 都为同构，就说这些边界条件满足\term[tuoyuan butiaojian]{椭圆补条件}[elliptic complementing condition]。各行分别按该边界算子的阶数取主部；在等维假设下，只需检验齐次边界条件没有非零衰减解。
\end{definition}

上一定义只检验非零切向频率。在一维区间上没有这样的方向，端点问题直接由常微分方程处理。在高维情形，补条件也不消除全局低频障碍：Neumann 问题仍可有常数核。\cref{g3:exm:boundary-dirichlet-robin-symbol}中的系统具有标量主符号 \(|\zeta|^2I\)，其衰减空间与维数已经直接算出，因而没有额外的根数假设需要验证。

\begin{example}\label{g3:exm:elliptic-bad-boundary-disk}
	在单位圆盘 \(D\subset\mathbb C\) 上，令 \(P=-\partial_x^2-\partial_y^2\)，并对复值函数施加
	\[
		Bu=(\partial_\nu+\mathrm i\partial_\theta)u=0
		\quad\text{于 }\partial D.
	\]
	这里 \(\theta\) 是圆周的弧长参数。每个 \(u_k(z)=z^k\) 都满足 \(Pu_k=0\)，且
	\[
		Bu_k=(k+\mathrm i\cdot\mathrm i k)\mathrm e^{\mathrm i k\theta}=0.
	\]
	因此齐次问题有无限维光滑核。展平边界后，最高阶边界算子为 \(-\partial_t+\mathrm i\partial_y\)，它对衰减模态的作用是 \((|\xi|-\xi)c\)，在 \(\xi>0\) 时恰好为零。这把全局失败与局部补条件的失败对应起来；内部算子仍是最简单的椭圆算子。
\end{example}

下一步采用一类可由能量和差商直接证明的系统：主符号为度量的平方范数乘恒等映射，边界条件由光滑的正交丛分裂给出。这个分裂发生在每个边界纤维内，允许不同分量承担不同条件；它不包含在边界的相邻区域之间突然切换条件的交界点。

\subsection{Hodge 的绝对与相对边界条件}
\label{g3:subsec:1403:03}
设 \((M,g)\) 为紧的光滑 \(n\) 维带边界黎曼流形，\(\nu\) 为外单位法向，\(j:\partial M\hookrightarrow M\) 为包含映射。可以使用实形式，也可以将形式复化；以下配对均在第一变量线性。对一个 \(q\) 次形式，在边界有正交分解
\begin{equation}\label{g3:eq:form-tangential-normal-traces}
	\omega|_{\partial M}=\tau+\nu^\flat\wedge\beta,
	\qquad \tau=\mathrm t\omega:=j^*\omega,
	\qquad \beta=\mathrm n\omega:=j^*(\iota_\nu\omega).
\end{equation}
右侧把边界上的形式视作其切向延拓。这里的 \(\mathrm n\omega\) 是 \(q-1\) 次形式；\(q=0\) 时规定它为零，\(q=n\) 时 \(\mathrm t\omega\) 自动为零。定义与积分均使用黎曼密度，不要求 \(M\) 可定向。

\begin{lemma}[形式的 Green 边界配对]\label{g3:lem:hodge-boundary-green}
	设 \(M\) 为紧的光滑带边界黎曼流形，\(\delta\) 为 \(\mathrm d\) 的内部形式伴随，\(\Delta=\mathrm d\delta+\delta\mathrm d\)。对直到边界都光滑的形式，有
	\begin{equation}\label{g3:eq:d-boundary-green}
		(\mathrm d\alpha,\eta)_M-(\alpha,\delta\eta)_M
		=(\mathrm t\alpha,\mathrm n\eta)_{\partial M}.
	\end{equation}
	若 \(\omega,\eta\) 同为 \(q\) 次形式，记
	\[
		Q(\omega,\eta)=(\mathrm d\omega,\mathrm d\eta)_M
		+(\delta\omega,\delta\eta)_M.
	\]
	则
	\begin{equation}\label{g3:eq:hodge-energy-green}
		Q(\omega,\eta)=(\Delta\omega,\eta)_M
		+(\mathrm n\mathrm d\omega,\mathrm t\eta)_{\partial M}
		-(\mathrm t\delta\omega,\mathrm n\eta)_{\partial M}.
	\end{equation}
	相应的二阶边界配对为
	\begin{align}\label{g3:eq:hodge-second-green}
		(\Delta\omega,\eta)_M-(\omega,\Delta\eta)_M
		={}&(\mathrm t\omega,\mathrm n\mathrm d\eta)_{\partial M}
		-(\mathrm n\mathrm d\omega,\mathrm t\eta)_{\partial M}\notag\\
		&+(\mathrm t\delta\omega,\mathrm n\eta)_{\partial M}
		-(\mathrm n\omega,\mathrm t\delta\eta)_{\partial M}.
	\end{align}
\end{lemma}
\begin{proof}
	在局部正交标架 \((e_i)\) 及其对偶标架 \((e^i)\) 中，无挠性给出
	\[
		\mathrm d=\sum_i e^i\wedge\nabla_{e_i},\qquad
		\delta=-\sum_i\iota_{e_i}\nabla_{e_i}.
	\]
	第二式可由第一式逐项分部积分得到；\(e^i\wedge\) 的纤维伴随是 \(\iota_{e_i}\)，相容联络及体积密度的导数合成散度项。对散度积分，边界项就是 \(\langle\nu^\flat\wedge\alpha,\eta\rangle=\langle\mathrm t\alpha,\mathrm n\eta\rangle\)，从而得到式\eqref{g3:eq:d-boundary-green}。将它分别用于 \((\omega,\mathrm d\eta)\) 与 \((\delta\omega,\eta)\)，并在需要时取复共轭，便得到式\eqref{g3:eq:hodge-energy-green}。再交换 \(\omega,\eta\) 并取共轭相减，得到式\eqref{g3:eq:hodge-second-green}。
\end{proof}

\begin{definition}\label{g3:def:absolute-relative-boundary}
	对紧光滑带边界黎曼流形上的光滑 \(q\) 次形式，下列两组条件分别称为 \term[juedui bianjie tiaojian]{绝对边界条件}[absolute boundary conditions]与\term[xiangdui bianjie tiaojian]{相对边界条件}[relative boundary conditions]：
	\begin{equation}\label{g3:eq:absolute-relative-boundary}
		\begin{array}{c|cc}
		 &\text{零阶条件}&\text{一阶条件}\\ \hline
		\mathrm A&\mathrm n\omega=0&\mathrm n\mathrm d\omega=0\\
		\mathrm R&\mathrm t\omega=0&\mathrm t\delta\omega=0.
		\end{array}
	\end{equation}
	它们的能量空间分别为
	\[
		V_{\mathrm A}^q=\{\omega\in H^1(M;\Lambda^q):\mathrm n\omega=0\},\qquad
		V_{\mathrm R}^q=\{\omega\in H^1(M;\Lambda^q):\mathrm t\omega=0\}.
	\]
	这里零阶条件按有界迹解释；一阶条件不放进 \(H^1\) 能量空间，而将由弱方程产生。
\end{definition}

对函数，绝对条件就是 Neumann 条件，相对条件就是 Dirichlet 条件。对一般次数的形式，两类条件使式\eqref{g3:eq:hodge-energy-green}的边界项消失，因此在相应光滑定义域上有 \((\Delta\omega,\omega)=Q(\omega,\omega)\ge0\)。但是，要从这个能量构造弱解，还必须证明它控制整个一阶 Sobolev 范数。

\begin{lemma}[边界形算子与一阶能量]\label{g3:lem:hodge-boundary-energy-estimate}
	设 \(M\) 为紧的光滑带边界黎曼流形，\(0\le q\le n\)。则存在 \(C\)，使 \(V_{\mathrm A}^q\) 或 \(V_{\mathrm R}^q\) 中的每个形式满足
	\begin{equation}\label{g3:eq:hodge-gaffney-estimate}
		\|\omega\|_{H^1}^2
		\le C\bigl(\|\mathrm d\omega\|_{L^2}^2
		+\|\delta\omega\|_{L^2}^2+\|\omega\|_{L^2}^2\bigr).
	\end{equation}
	更具体地，令 \(S(X)=\nabla_X\nu\) 作用于边界切向量，并定义它在切向 \(p\) 次形式上的作用
	\[
		S^{[p]}=\sum_{a=1}^{n-1}e^a\wedge\iota_{S e_a},
		\qquad S^{[0]}=0.
	\]
	存在光滑自伴零阶丛映射 \(\mathcal R_q\)，使光滑形式分别满足
	\begin{align}
		Q(\omega,\omega)
		&=\|\nabla\omega\|_{L^2}^2+(\mathcal R_q\omega,\omega)_M
		+(S^{[q]}\mathrm t\omega,\mathrm t\omega)_{\partial M},
		&&\mathrm n\omega=0,\label{g3:eq:hodge-energy-absolute}\\
		Q(\omega,\omega)
		&=\|\nabla\omega\|_{L^2}^2+(\mathcal R_q\omega,\omega)_M
		+\bigl(((\operatorname{tr}S)I-S^{[q-1]})\mathrm n\omega,
		\mathrm n\omega\bigr)_{\partial M},
		&&\mathrm t\omega=0.\label{g3:eq:hodge-energy-relative}
	\end{align}
	在没有相应分量的端点次数，忽略该边界项。这里 \(S=\nabla\nu\) 的符号约定与本卷采用的形算子 \(A_\nu=-\nabla\nu\) 相反。
\end{lemma}
\begin{proof}
	先说明内部二阶项怎样合并。写 \(\varepsilon_i\omega=e^i\wedge\omega\)、\(\iota_i=\iota_{e_i}\)。外乘与收缩满足
	\(\varepsilon_i\iota_j+\iota_j\varepsilon_i=\delta_{ij}I\)。在一点选取协变导数为零的正交标架，将前一引理中的 \(\mathrm d,\delta\) 展开，得到
	\begin{equation}\label{g3:eq:hodge-rough-laplacian}
		\Delta=\nabla^*\nabla+\mathcal R_q,
		\qquad
		\mathcal R_q=-\sum_{i,j}\varepsilon_i\iota_j R^\Lambda(e_i,e_j).
	\end{equation}
	这里 \(R^\Lambda(X,Y)=[\nabla_X,\nabla_Y]-\nabla_{[X,Y]}\) 是形式丛上的曲率。确切地说，两种二阶导数之和先成为 \(-\sum_i\nabla_{e_i}\nabla_{e_i}\)，剩余的反对称部分正是所写交换子；恢复非正规标架时，迹联络项组成 \(\nabla^*\nabla\)。因此剩余项是光滑零阶映射。\(\Delta\) 与 \(\nabla^*\nabla\) 在内部紧支截面上均形式自伴，故其零阶差逐点自伴。紧性使 \(\mathcal R_q\) 的范数有界。

	现在计算边界。沿法向测地线延拓 \(\nu\) 及边界切向正交标架，使 \(\nabla_\nu\nu=0\)、\(\nabla_\nu e_a=0\)。\(S\) 自伴，因为对切向 \(X,Y\)，无挠性给出
	\(\langle SX,Y\rangle-\langle X,SY\rangle=-\langle\nu,[X,Y]\rangle=0\)。当 \(\mathrm n\omega=0\) 时，沿边界求导得到
	\[
		\iota_\nu\nabla_{e_a}\omega=-\iota_{S e_a}\omega,
	\]
	因而
	\begin{equation}\label{g3:eq:absolute-robin-form}
		\mathrm n\mathrm d\omega
		=(\nabla_\nu\omega)_{\mathrm T}+S^{[q]}\mathrm t\omega.
	\end{equation}
	这里下标 \(\mathrm T\) 表示切向分量。确实，展开 \(\iota_\nu\mathrm d\omega\)，法向项给出 \((\nabla_\nu\omega)_{\mathrm T}\)，切向项为
	\(-\sum_a e^a\wedge\iota_\nu\nabla_{e_a}\omega\)，代入上一式即得结论。

	当 \(\mathrm t\omega=0\) 时，边界上写 \(\omega=\nu^\flat\wedge\beta\)。\(\delta\) 的法向项给出 \(-\nabla_\nu\beta\)。切向项的切向部分为
	\[
		-\sum_a\iota_{e_a}\bigl((S e_a)^\flat\wedge\beta\bigr)
		=-(\operatorname{tr}S)\beta+S^{[q-1]}\beta.
	\]
	含有 \(\nu^\flat\) 的其他项取切向分量后为零。因此
	\begin{equation}\label{g3:eq:relative-robin-form}
		\mathrm t\delta\omega
		=-\bigl(\mathrm n\nabla_\nu\omega
		+((\operatorname{tr}S)I-S^{[q-1]})\mathrm n\omega\bigr).
	\end{equation}
	这两式也说明绝对、相对条件分别是一个分量的 Dirichlet 条件，加上互补分量的 Robin 条件。

	粗 Laplace 算子的分部积分为
	\[
		(\nabla^*\nabla\omega,\omega)_M
		=\|\nabla\omega\|_{L^2}^2-(\nabla_\nu\omega,\omega)_{\partial M}.
	\]
	把它与式\eqref{g3:eq:hodge-energy-green}、\eqref{g3:eq:hodge-rough-laplacian}合并。在绝对零阶条件下，用式\eqref{g3:eq:absolute-robin-form}，法向导数项相消，留下式\eqref{g3:eq:hodge-energy-absolute}；相对情形用式\eqref{g3:eq:relative-robin-form}，同样得到式\eqref{g3:eq:hodge-energy-relative}。\(S^{[p]}\) 自伴：在 \(S\) 的正交特征标架中，基形式 \(e^{a_1}\wedge\cdots\wedge e^{a_p}\) 的特征值是对应 \(p\) 个特征值之和。因此两种边界零阶项均为实数且绝对值不超过 \(C\|\operatorname{Tr}\omega\|_{L^2(\partial M)}^2\)。

	由丛值迹估计\cref{g3:prop:sobolev-boundary-bundles}，对每个 \(\varepsilon>0\)，
	\[
		\|\operatorname{Tr}\omega\|_{L^2(\partial M)}^2
		\le\varepsilon\|\nabla\omega\|_{L^2}^2+C_\varepsilon\|\omega\|_{L^2}^2.
	\]
	先取 \(\varepsilon\) 足够小，吸收梯度项，便得到光滑形式的式\eqref{g3:eq:hodge-gaffney-estimate}。为延拓到所述能量空间，在边界适配标架中把要求零迹的分量作奇反射，其他分量作偶反射，再用关于法向偶对称的光滑核卷积。奇分量仍在边界为零；局部截断、单位分解及内部卷积使这些光滑截面在 \(H^1\) 中逼近原截面。所有项对 \(H^1\) 连续，故估计以及相应能量恒等式都延续到 \(V_{\mathrm A}^q,V_{\mathrm R}^q\)。
\end{proof}

式\eqref{g3:eq:hodge-gaffney-estimate}保留 \(L^2\) 项，是因为存在的调和形式可能有零能量。它已足以使 \(Q+(\cdot,\cdot)_{L^2}\) 在两种能量空间上强制。剩下的问题是：这样得到的弱解是否具有二阶乃至更高的正则性，以及它的一阶边界条件是否真正成立。下面把这一步放在统一的丛值问题中证明。

\subsection{Laplace 型算子的混合边值问题}
\label{g3:subsec:1403:04}
设 \(E\to M\) 为光滑 Hermitian 向量丛，\(\nabla\) 为相容联络。称
\[
	P=\nabla^*\nabla+R
\]
为这里使用的 Laplace 型算子，其中 \(R=R^*\) 是光滑丛映射。它的主符号是 \(|\xi|_g^2I\)。在边界给定光滑正交分解
\[
	E|_{\partial M}=E_{\mathrm D}\oplus E_{\mathrm N},
\]
其秩在每个边界连通分支上固定，并给定 \(E_{\mathrm N}\) 上光滑自伴映射 \(S\)。本小节的 \(S\) 表示任意 Robin 系数；应用于形式时再取上一小节计算出的相应映射。定义
\begin{align}
	V&=\{u\in H^1(M;E):\Pi_{\mathrm D}\operatorname{Tr}u=0\},\label{g3:eq:mixed-bundle-energy-space}\\
	a(u,v)&=(\nabla u,\nabla v)_M+(Ru,v)_M
	+(S\Pi_{\mathrm N}u,\Pi_{\mathrm N}v)_{\partial M}.
	\label{g3:eq:mixed-bundle-energy-form}
\end{align}
边界积分中的截面均取迹。对光滑 \(u\)，分部积分表明 \(a(u,v)=(Pu,v)\) 对所有 \(v\in V\) 成立的条件是
\begin{equation}\label{g3:eq:mixed-bundle-boundary}
	\Pi_{\mathrm D}u=0,\qquad
	\Pi_{\mathrm N}\nabla_\nu u+S\Pi_{\mathrm N}u=0.
\end{equation}
在边界适配标架中，两个投影成为常数矩阵；沿法向平行延拓标架后，法向联络导数就是普通分量导数。主符号及边界主部因此正好是\cref{g3:exm:boundary-dirichlet-robin-symbol}。

\begin{lemma}[固定分量分裂的边界正则性]\label{g3:lem:mixed-bundle-local-regularity}
	设 \(Q_R^+=(-R,R)^{n-1}\times(0,R)\)，\(\Gamma_R=\{t=0\}\cap\partial Q_R^+\)，\(\mathbb C^r=E_{\mathrm D}\oplus E_{\mathrm N}\) 为固定正交分解。设 \(a^{ij}\) 为实对称光滑一致正定矩阵，\(S\) 为 \(E_{\mathrm N}\) 上的光滑矩阵。若 \(u\in H^1(Q_R^+;\mathbb C^r)\)、\(\Pi_{\mathrm D}\operatorname{Tr}u=0\)，且对侧面及顶面附近消失、满足相同零阶边界条件的全部 \(v\in H^1\)，有
	\begin{equation}\label{g3:eq:mixed-flat-weak-equation}
		\int_{Q_R^+}a^{ij}\langle\partial_j u,\partial_i v\rangle
		+\int_{\Gamma_R}\langle Su_{\mathrm N},v_{\mathrm N}\rangle
		=\int_{Q_R^+}\langle F,v\rangle
		+\int_{Q_R^+}\langle G^i,\partial_i v\rangle
		+\int_{\Gamma_R}\langle g,v_{\mathrm N}\rangle,
	\end{equation}
	其中 \(F\in L^2\)、\(G\in H^1\)、\(g\in H^{1/2}(\Gamma_R;E_{\mathrm N})\)，则对 \(0<r_0<R\)，有 \(u\in H^2(Q_{r_0}^+)\) 及
	\begin{equation}\label{g3:eq:mixed-flat-h2-estimate}
		\|u\|_{H^2(Q_{r_0}^+)}
		\le C\bigl(\|F\|_{L^2}+\|G\|_{H^1}
		+\|g\|_{H^{1/2}}+\|u\|_{H^1}\bigr).
	\end{equation}
	局部边界范数以图中的半整数 Sobolev 范数解释。向弱形式加入光滑矩阵系数的内部一阶或零阶项，结论仍成立，常数还依赖这些系数的有界范数。
\end{lemma}
\begin{proof}
	先消去独立的边界源项 \(g\)。在稍大的边界片上乘以切向截断，使 \(g\) 在所需较小片上不变。由\cref{g3:thm:half-integer-boundary-trace}的联合迹右逆，存在 \(w\in H^2\)，满足全分量的 \(\operatorname{Tr}w=0\) 及
	\[
		-a^{nn}\partial_t w=g\quad\text{于边界},
	\]
	右端视为 \(E_{\mathrm N}\) 中的数据，且 \(\|w\|_{H^2}\le C\|g\|_{H^{1/2}}\)。这是因为 \(a^{nn}\) 有正下界，而零迹保证所有切向导数迹为零，故外共法向导数正是 \(-a^{nn}\partial_t w\)。把 \(u-w\) 代入弱式，并对 \(w\) 分部积分，消去 \(g\)；新增内部源项是 \(L^2\) 函数，由 \(\|w\|_{H^2}\) 控制。局部截断只需在严格较大的片中进行，故不影响欲证的较小半柱估计。以下取 \(g=0\)。

	沿一个切向方向取差商 \(\delta_h\)，令 \(z=\delta_hu\)，并以
	\[
		v=-\delta_{-h}(\eta^2z)
	\]
	测试，其中 \(\eta\) 在较小半柱上等于一，在人工侧面及顶面附近消失。由于 \(\Pi_{\mathrm D}\) 是常数矩阵，切向平移保持其零迹，所以这个测试函数可用。主部、截断、\(F\) 和 \(G\) 的计算逐分量同\cref{g3:lem:boundary-tangential-differences}，并对有限个分量求和：主项至少为 \(c\|\eta\mathrm Dz\|_2^2\)，其余内部项至多为
	\[
		\varepsilon\|\eta\mathrm Dz\|_2^2
		+C_\varepsilon(\|F\|_2+\|G\|_{H^1}+\|u\|_{H^1})^2.
	\]
	在此仍把 \(F\) 上的差商留在测试函数，利用
	\(\|\delta_{-h}(\eta^2z)\|_2\le C(\|\eta\mathrm Dz\|_2+\|z\|_2)\)，所以没有要求 \(F\) 的导数。

	新增的边界项经切向差商分部积分成为
	\[
		\int_{\Gamma_R}
		\left\langle
		S(y+he_a)z_{\mathrm N}+(\delta_hS)u_{\mathrm N},
		\eta^2z_{\mathrm N}\right\rangle.
	\]
	其绝对值由
	\[
		C\|\operatorname{Tr}(\eta z)\|_2^2
		+C\|\operatorname{Tr}u\|_2^2
	\]
	控制。对 \(\eta z\) 使用带任意小系数的迹估计，再用
	\(\|z\|_2\le C\|u\|_{H^1}\)，可将第一项界为
	\(\varepsilon\|\eta\mathrm Dz\|_2^2+C_\varepsilon\|u\|_{H^1}^2\)。第二项直接由有界迹控制。选取 \(\varepsilon\) 后吸收，得到与 \(h\) 无关的 \(\|\eta\mathrm D\delta_hu\|_2\) 界。差商弱收敛判据遂给出所有含切向指标的二阶导数。

	对内部紧支测试函数，边界项消失，弱式给出
	\[
		\partial_n(a^{nn}\partial_nu)
		=-F+\partial_iG^i
		-\sum_{(i,j)\ne(n,n)}\partial_i(a^{ij}\partial_ju).
	\]
	右侧各分量均在 \(L^2\)，因此先求通量的弱导数，再乘 \((a^{nn})^{-1}\)，便得 \(\partial_n^2u\in L^2\) 及所列估计。一维情形直接使用这一步。最后，内部低阶项
	\(\langle B^j\partial_ju+Eu,v\rangle+\langle C^iu,\partial_iv\rangle\)
	可分别移入 \(F\) 与 \(G\)；光滑系数使 \(B^j\partial_ju+Eu\in L^2\)、\(C^iu\in H^1\)，其范数都由 \(\|u\|_{H^1}\) 控制。这证明了矩阵低阶项的版本。
\end{proof}

\begin{theorem}[光滑丛分裂下的 Fredholm 边值问题]\label{g3:thm:mixed-bundle-fredholm}
	设 \(M\) 为紧的光滑 \(n\ge1\) 维带边界黎曼流形，\(E\to M\) 为有限正秩光滑 Hermitian 向量丛，\(\nabla\) 为相容联络，\(R=R^*\) 为光滑丛映射。设边界上的光滑正交分解 \(E_{\mathrm D}\oplus E_{\mathrm N}\) 及光滑自伴 \(S:E_{\mathrm N}\to E_{\mathrm N}\) 如上，令 \(P=\nabla^*\nabla+R\)。对整数 \(k\ge0\)，定义闭子空间
	\[
		H_B^{k+2}=
		\{u\in H^{k+2}(M;E):
		\Pi_{\mathrm D}u=0,\quad
		\Pi_{\mathrm N}\nabla_\nu u+S\Pi_{\mathrm N}u=0\}.
	\]
	边界条件按 \(H^{k+3/2}\) 与 \(H^{k+1/2}\) 迹解释。则有以下结论。
	\begin{enumerate}
		\item 若 \(u\in V\) 满足 \(a(u,v)=(f,v)\) 对所有 \(v\in V\) 成立，且 \(f\in H^k(M;E)\)，则 \(u\in H_B^{k+2}\)，并有
		\begin{equation}\label{g3:eq:mixed-bundle-global-estimate}
			\|u\|_{H^{k+2}}
			\le C_k(\|f\|_{H^k}+\|u\|_{L^2}).
		\end{equation}
		\item 在 \(L^2(M;E)\) 上取定义域 \(\mathcal D(A_B)=H_B^2\)，令 \(A_Bu=Pu\)。则 \(A_B\) 自伴且下有界；充分大的 \(\lambda>0\) 使 \((A_B+\lambda I)^{-1}\) 存在且紧。其特征截面可取为直到边界光滑的正交归一基，特征值为实数，重数有限，按重数排列后趋于 \(+\infty\)。
		\item \(P:H_B^{k+2}\to H^k(M;E)\) 是指标为零的 Fredholm 算子。其核 \(\mathcal H_B\) 是同一个有限维光滑截面空间，而
		\begin{equation}\label{g3:eq:mixed-bundle-fredholm-range}
			\operatorname{ran}P
			=\{f\in H^k:(f,h)_{L^2}=0\text{ 对所有 }h\in\mathcal H_B\}.
		\end{equation}
		对上述右端存在唯一的 \(u\in H_B^{k+2}\cap\mathcal H_B^\perp\)，且
		\(\|u\|_{H^{k+2}}\le C_k\|f\|_{H^k}\)。
	\end{enumerate}
\end{theorem}
\begin{proof}
	证明先从能量得到一个可逆平移，再用边界正则性辨认其定义域，最后由紧谱定理除去平移。由于 \(R,S\) 有界，迹的小系数估计给出
	\[
		a(u,u)\ge\tfrac12\|\nabla u\|_2^2-C\|u\|_2^2,
		\qquad u\in V.
	\]
	型 \(a\) 在 \(V\times V\) 上连续且 Hermitian。\(V\) 是 \(H^1\) 的闭子空间；取 \(\lambda\) 足够大后，\(a_\lambda=a+\lambda(\cdot,\cdot)\) 强制。由\cref{g3:thm:lax-milgram}，每个 \(f\in L^2\) 对应唯一 \(u\in V\)，满足
	\(a_\lambda(u,v)=(f,v)\)，且 \(\|u\|_{H^1}\le C\|f\|_2\)。

	对任一满足 \(a(u,v)=(f,v)\) 的弱解，取边界法向坐标与适配丛标架。由相容联络，主部能量逐分量为 \(a^{ij}\partial_j u\,\overline{\partial_i v}\)，其中 \(a^{ij}=\rho g^{ij}\) 为实对称正定标量矩阵；\(\rho\) 是体积密度。联络系数只产生前一引理允许的矩阵低阶项，边界密度并入 \(S\)。两个边界投影为常数矩阵。因此该引理给出局部 \(H^2\) 正则性。内部同样逐分量使用差商，矩阵低阶项移入源项，得到内部 \(H^2\) 界。有限覆盖得到
	\[
		\|u\|_{H^2}\le C(\|f\|_2+\|u\|_{H^1}).
	\]
	此时对 \(u\) 分部积分合法，比较弱等式得
	\[
		(\Pi_{\mathrm N}\nabla_\nu u+S\Pi_{\mathrm N}u,
		\Pi_{\mathrm N}\operatorname{Tr}v)_{\partial M}=0
		\quad(v\in V).
	\]
	每个光滑 \(E_{\mathrm N}\) 边界截面都有沿领圈延拓再截断得到的 \(V\) 中截面；这些边界截面稠密。因此一阶边界迹为零，弱解确实属于 \(H_B^2\)。自然边界条件在这一步由方程得出。

	下面求高阶导数。假定已知 \(f\in H^k\)、\(u\in H^{k+1}\)，其中 \(k\ge1\)。在稍小边界片内，对弱形式求 \(k\) 次切向导数，令 \(w=\partial_{y'}^\beta u\)、\(|\beta|=k\)。常数 Dirichlet 投影使 \(\Pi_{\mathrm D}\operatorname{Tr}w=0\)。主部交换子在右端产生
	\[
		G_\beta^i=-\sum_{\gamma<\beta}\binom\beta\gamma
		(\partial_{y'}^{\beta-\gamma}a^{ij})
		\partial_j\partial_{y'}^\gamma u,
	\]
	其中 \(\gamma<\beta\) 表示逐分量不超过 \(\beta\) 且不等于 \(\beta\)。每项属于 \(H^1\)，因为 \(|\gamma|\le k-1\) 且已知 \(u\in H^{k+1}\)。内部低阶交换子或归入 \(L^2\) 的 \(F_\beta\)，或归入同样的 \(H^1\) 向量源项；对 \(f\) 最多求 \(k\) 次导数。边界项则给出
	\[
		g_\beta=-\sum_{\gamma<\beta}\binom\beta\gamma
		(\partial_{y'}^{\beta-\gamma}S)
		\partial_{y'}^\gamma u_{\mathrm N}\in H^{1/2}.
	\]
	最后的包含由半整数迹定理及光滑乘法得到：\(\partial_{y'}^\gamma u\) 至少属于 \(H^2\)，其边界迹甚至属于 \(H^{3/2}\)。这些等式可先在分布意义下对测试函数求导，再用已有 \(H^{k+1}\) 正则性解释各项。局部截断带来的项具有相同的低阶形式。

	对 \(w\) 使用\cref{g3:lem:mixed-bundle-local-regularity}，获得所有含至少 \(k\) 个切向指标的 \(k+2\) 阶导数。为取得其余导数，把内部方程写成
	\[
		-a^{nn}\partial_n^2u
		=f_0+\sum_{(i,j)\ne(n,n)}a^{ij}\partial_i\partial_ju
		+\sum_j B^j\partial_ju+Cu,
	\]
	其中 \(f_0\) 与原来的 \(f\) 相差光滑密度因子。对总阶数为 \(k+2\) 的导数按法向指标数递增归纳。已知法向指标数至多二的情形；若目标有 \(r\ge3\) 个法向指标，就对上式求相应的 \(k\) 阶导数。右侧的最高阶导数至多有 \(r-1\) 个法向指标，已在此前得到；落在系数上的导数使 \(u\) 的阶数降至 \(k+1\) 以下。除以正的标量 \(a^{nn}\)，便恢复目标导数。内部图用各方向的同类差商与求导。有限覆盖给出
	\[
		\|u\|_{H^{k+2}}\le C_k(\|f\|_{H^k}+\|u\|_{H^{k+1}}).
	\]
	一维情形全部由法向方程递推。\(H^m\) 的延拓、Fourier 插值及有限分解给出
	\(\|u\|_{H^{k+1}}\le\varepsilon\|u\|_{H^{k+2}}+C_\varepsilon\|u\|_2\)。取小 \(\varepsilon\) 吸收；\(k=0\) 时同样处理 \(H^1\) 项，得到式\eqref{g3:eq:mixed-bundle-global-estimate}。这也证明光滑右端的弱解直到边界光滑。

	回到强制型得到的平移方程。上述正则性对 \(P+\lambda I\) 同样适用，所以其弱解属于 \(H_B^2\)。反过来，\(H_B^2\) 中的截面由分部积分满足相同弱等式。因此
	\[
		K=(A_B+\lambda I)^{-1}:L^2\longrightarrow H_B^2
	\]
	是有界映射，并且由 \(H^2\hookrightarrow L^2\) 的紧嵌入成为紧算子。若 \(u=Kf,v=Kg\)，则
	\[
		(f,Kg)=a_\lambda(Kf,Kg)=(Kf,g),\qquad
		(Kf,f)=a_\lambda(Kf,Kf)>0\quad(f\ne0).
	\]
	所以 \(K\) 自伴、非负且单射。对称性已经由边界配对保证；为核对无界算子的真正自伴性，若 \(v\in\mathcal D(A_B^*)\)、\(A_B^*v=w\)，则对全部 \(u\in H_B^2\)，
	\[
		((A_B+\lambda I)u,v)=(u,w+\lambda v)
		=((A_B+\lambda I)u,K(w+\lambda v)).
	\]
	平移算子满射，故 \(v=K(w+\lambda v)\in H_B^2\)，证明伴随定义域没有扩大。定义域稠密则来自内部紧支光滑截面包含在其中。

	\cref{g3:thm:compact-positive-spectrum}给出 \(Ke_j=\mu_je_j\)，其中 \((e_j)\) 为 \(L^2\) 正交归一基，\(\mu_j>0\)、\(\mu_j\to0\)。于是
	\[
		A_Be_j=\lambda_je_j,\qquad
		\lambda_j=\mu_j^{-1}-\lambda\longrightarrow+\infty.
	\]
	正则性递推使每个 \(e_j\) 光滑；有限重数来自 \(K\) 的紧性。由于 \(H_B^2=\operatorname{ran}K\)，一个展开 \(u=\sum_j u_je_j\) 属于定义域当且仅当
	\(\sum_j\mu_j^{-2}|u_j|^2<\infty\)。这给出该实现的全部谱，并说明零特征值只有有限重数。

	记 \(\mathcal H_B=\ker A_B\)。对 \(f\perp\mathcal H_B\)，级数
	\[
		u=\sum_{\lambda_j\ne0}\frac{(f,e_j)}{\lambda_j}e_j
	\]
	属于 \(H_B^2\)：非零 \(\lambda_j\) 离零有正距离，而且 \(\mu_j^{-1}/|\lambda_j|\) 一致有界。它是唯一正交于核的解，并有 \(\|u\|_2\le C\|f\|_2\)。若 \(f\in H^k\)，第一部分再给出 \(u\in H_B^{k+2}\) 及加强估计。必要的正交条件由自伴性得到，故式\eqref{g3:eq:mixed-bundle-fredholm-range}成立。其右侧是 \(H^k\) 中的闭子空间，余维恰为 \(\dim\mathcal H_B\)：核中的光滑正交基可直接检验这些泛函线性独立。核与余核同维，指标为零。
\end{proof}

\begin{corollary}[形式的绝对与相对实现]\label{g3:cor:hodge-boundary-fredholm}
	设 \(M\) 为紧的光滑 \(n\ge1\) 维带边界黎曼流形，\(0\le q\le n\)。\(\Delta_q=\mathrm d\delta+\delta\mathrm d\) 在 \(L^2(M;\Lambda^q)\) 上分别取绝对或相对 \(H^2\) 定义域，得到非负自伴算子 \(\Delta_{q,\mathrm A}\)、\(\Delta_{q,\mathrm R}\)。二者都有紧预解式、光滑特征形式与有限维核，而且
	\[
		\mathcal H_{\mathrm A}^q
		=\{\omega\in V_{\mathrm A}^q:\mathrm d\omega=0,\ \delta\omega=0\},
		\qquad
		\mathcal H_{\mathrm R}^q
		=\{\omega\in V_{\mathrm R}^q:\mathrm d\omega=0,\ \delta\omega=0\}.
	\]
	右侧最初按弱导数解释，其中每个元素实际上直到边界光滑。对整数 \(k\ge0\)，右端 \(f\in H^k\) 可解相应齐次边值问题，当且仅当 \(f\) 正交于相应调和核；正交于核的解唯一且满足 \(\|\omega\|_{H^{k+2}}\le C_k\|f\|_{H^k}\)。
\end{corollary}
\begin{proof}
	由式\eqref{g3:eq:hodge-rough-laplacian}，取 \(E=\Lambda^q\)、\(R=\mathcal R_q\)。绝对情形令法向形式为 Dirichlet 分量、切向形式为互补分量，并取 \(S=S^{[q]}\)；相对情形交换这两个分量，在 \(\nu^\flat\wedge\Lambda^{q-1}T^*\partial M\) 上取 \(S=(\operatorname{tr}S_0)I-S_0^{[q-1]}\)，其中 \(S_0(X)=\nabla_X\nu\)。式\eqref{g3:eq:absolute-robin-form}、\eqref{g3:eq:relative-robin-form}说明两种丛边界条件正是所定义的绝对与相对条件，所需系数光滑自伴。

	两种能量恒等式及极化表明相应的 \(a\) 等于 \(Q\)，故非负性成立；式\eqref{g3:eq:hodge-gaffney-estimate}还保证 \(Q+(\cdot,\cdot)\) 强制，可以直接取平移 \(\lambda=1\)。其余正则性、谱与 Fredholm 结论来自上一定理。若 \(\omega\) 属于算子核，能量为零，故 \(\mathrm d\omega=\delta\omega=0\)。反之，若 \(\omega\in V_{\mathrm A}^q\) 或 \(V_{\mathrm R}^q\) 且两个弱导数均为零，则 \(Q(\omega,v)=0\) 对相应全部测试形式成立。上一定理将它提升为光滑齐次边值问题的解，故它属于算子核。
\end{proof}

这里的证明用到了三项具体结构：标量的正定主符号、边界上光滑的纤维分裂，以及边界项只有零阶的 Robin 矩阵。一般系统的补条件理论还允许矩阵主符号、不同阶数的边界算子及非局部边界条件；其证明需要构造与衰减空间相容的边界参数解算子。若引入谱参数，还须对参数与切向频率共同趋大的情形检验相应条件。半空间检验给出这类理论的局部问题，不能单凭一次冻结计算就代替全局正则性与 Fredholm 证明。
% 纲要条目排为编号小节标题，并进入目录与书签；说明段仍保留于纲要注释中。

% !TeX root = ../../Book3.tex
% 纲要标题：### 14.4 算子族的参数依赖与核维数跳跃
% 纲要位置：几何学计划纲要.md，第 2695 行。
% 正文状态：初稿；尚未正式审核。

% 纲要内容开始
% 
% * 固定核维数下的投影和 Green 算子
% * 参数椭圆估计和连续性
% * 核维数跳跃及谱流动机
% * 参数化求解与几何族
% 
% ---
% 
% 纲要内容结束

\section{算子族的参数依赖与核维数跳跃}
\label{g3:sec:1404}

方程的系数、度量或区域发生变化时，我们希望解也随之连续变化。这里首先要说明“变化”在哪个空间中度量。例如，Dirichlet 条件可以给出固定的定义域，而随度量改变的法向导数通常不能。以下先在固定 Banach 空间之间研究算子族，再把边值问题放进这个模型。第十三章的有限维消元将再次出现：可逆的大块控制求解，有限维的小块记录核的变化。

\subsection{固定核维数下的投影和 Green 算子}
\label{g3:subsec:1404:01}
设 \(U\subset\mathbb R^s\) 为开集，\(X,Y\) 为固定的实或复 Banach 空间。说 \(a\mapsto T_a\) 为 \(C^r\) 算子族，指它作为取值于 \(\mathcal L(X,Y)\) 的映射为 \(C^r\)；\(r=0\) 表示算子范数连续。参数始终取实变量，即使 \(X,Y\) 是复空间。这一定义同时固定了定义域、值域及其范数。

\begin{theorem}\label{g3:thm:parameter-fredholm-constant-kernel}
	设 \(X,Y\) 为 Banach 空间，\(U\subset\mathbb R^s\) 开，\(a_0\in U\)，\(T:U\to\mathcal L(X,Y)\) 为 \(C^r\)，\(0\le r\le\infty\)。假设 \(T_{a_0}\) 为 Fredholm 算子，且在 \(a_0\) 附近 \(\dim\ker T_a=\dim\ker T_{a_0}\)。则在一个邻域内，核、值域分别有 \(C^r\) 的有界投影 \(P_a,Q_a\)，并有 \(C^r\) 映射 \(\widetilde G_a:Y\to X\)，满足
	\[
		T_a\widetilde G_a=Q_a,\qquad
		\widetilde G_aT_a=I-P_a.
	\]
	各 \(T_a\) 均为 Fredholm 算子，余核维数也恒定。这里的投影由选定补空间构造，未要求正交。
\end{theorem}
\begin{proof}
	由\cref{g3:prop:banach-fredholm-splitting}，取连续直和
	\[
		X=X_1\oplus K,\qquad Y=R\oplus C,
		\quad K=\ker T_{a_0},\quad R=\operatorname{ran}T_{a_0},
	\]
	其中 \(K,C\) 有限维，\(A_0=T_{a_0}|_{X_1}:X_1\to R\) 有有界逆。按此顺序写
	\[
		T_a=\begin{pmatrix}A_a&B_a\\ C_a&D_a\end{pmatrix},
		\qquad S_a=D_a-C_aA_a^{-1}B_a:K\to C.
	\]
	在 \(a_0\) 的小邻域，\(A_a\) 可逆。事实上当
	\(\|A_0^{-1}(A_a-A_0)\|<1\) 时，Neumann 级数收敛且给出其逆。逆映射为光滑映射：恒等式
	\[
		(A+H)^{-1}-A^{-1}=-(A+H)^{-1}HA^{-1}
	\]
	先给连续性，再给导数 \(H\mapsto-A^{-1}HA^{-1}\)；反复微分得到全部高阶导数。因此以下各个含 \(A_a^{-1}\) 的表达式均为 \(C^r\)。

	与式\eqref{g3:eq:fredholm-block-reduction}相同的行、列消元将 \(T_a\) 化为 \(\operatorname{diag}(A_a,S_a)\)。特别地
	\[
		\ker T_a=\{(-A_a^{-1}B_ak,k):k\in\ker S_a\}.
	\]
	恒核维数假设意味着 \(\dim\ker S_a=\dim K\)，故 \(S_a=0\)。记
	\[
		J_ak=(-A_a^{-1}B_ak,k),\qquad E_a=C_aA_a^{-1}:R\to C.
	\]
	于是 \(\ker T_a=J_aK\)，而 \(\operatorname{ran}T_a\) 恰为闭图
	\(\{(y,E_ay):y\in R\}\)。它以 \(\{0\}\oplus C\) 为补空间，故余维恒为 \(\dim C\)。所求算子可直接取为
	\begin{equation}\label{g3:eq:parameter-oblique-green}
		P_a(x,k)=J_ak,\qquad
		Q_a(y,c)=(y,E_ay),\qquad
		\widetilde G_a(y,c)=(A_a^{-1}y,0).
	\end{equation}
	\(P_a^2=P_a\)、\(Q_a^2=Q_a\) 由公式立即得到。又因 \(D_a=E_aB_a\)，直接相乘得到两个所述恒等式。全部映射有界，并按算子范数为 \(C^r\)。
\end{proof}

对椭圆方程，规范解通常用 \(L^2\) 正交条件选出。斜投影有利于计算，正交投影则能去掉补空间的选择。下面用有限个向量的 Gram 矩阵完成这一步。

\begin{theorem}\label{g3:thm:parameter-orthogonal-green}
	设 \(H_0,H_1\) 为实或复 Hilbert 空间，Banach 空间 \(X\) 连续且稠密地嵌入 \(H_0\)。设 \(U\subset\mathbb R^s\) 开，\(a\mapsto T_a\in\mathcal L(X,H_1)\) 为 \(C^r\) Fredholm 族，\(0\le r\le\infty\)，且核维数局部恒定。则 \(H_0\) 中的核正交投影 \(\Pi_a\)、\(H_1\) 中的值域正交补投影 \(\Pi_a^{\mathrm c}\)，分别按
	\[
		a\mapsto\Pi_a\in\mathcal L(H_0,X),\qquad
		a\mapsto\Pi_a^{\mathrm c}\in\mathcal L(H_1,H_1)
	\]
	为 \(C^r\)。存在唯一的 \(C^r\) 族 \(G_a\in\mathcal L(H_1,X)\)，满足
	\begin{equation}\label{g3:eq:parameter-orthogonal-green}
		T_aG_a=I-\Pi_a^{\mathrm c},\quad
		G_aT_a=I-\Pi_a,\quad
		\Pi_aG_a=0,\quad G_a\Pi_a^{\mathrm c}=0.
	\end{equation}
	把 \(T_a\) 看成从 \(H_0\) 到 \(H_1\)、域为 \(X\) 的稠定算子 \(\mathsf T_a\)，则它是闭算子，且
	\(\operatorname{ran}\Pi_a^{\mathrm c}=\ker\mathsf T_a^*\)。这里的伴随针对 \(H_0,H_1\) 的内积。
\end{theorem}
\begin{proof}
	固定 \(a_0\)，在上一证明中取 \(C=R^\perp\subset H_1\)。给有限维 \(K\) 固定一个内积，将 \(J_a:K\to X\hookrightarrow H_0\) 看成单射。其 Gram 算子 \(J_a^*J_a:K\to K\) 正定可逆，故
	\[
		\Pi_a=J_a(J_a^*J_a)^{-1}J_a^*.
	\]
	这个算子在 \(J_aK\) 上为恒等，且对全部 \(k\in K\)，\(u-\Pi_au\) 与 \(J_ak\) 正交，因此确为所求正交投影。有限维 Gram 算子及其逆均为 \(C^r\)，而 \(J_a\) 按 \(K\to X\) 的范数为 \(C^r\)，所以投影甚至按 \(H_0\to X\) 的范数为 \(C^r\)。

	值域是 \(E_a:R\to C\) 的图，其正交补是单射
	\[
		Z_a:C\to H_1,\qquad Z_ac=(-E_a^*c,c)
	\]
	的像。因为 \(\langle(y,E_ay),(-E_a^*c,c)\rangle=0\)，且每个与图正交的 \((z,c)\) 都满足 \(z=-E_a^*c\)，这里没有遗漏任何正交向量。故
	\[
		\Pi_a^{\mathrm c}=Z_a(Z_a^*Z_a)^{-1}Z_a^*
	\]
	也是 \(C^r\)。零维的 \(K\) 或 \(C\) 对应零投影。

	取上一结果的 \(\widetilde G_a\)，令
	\[
		G_a=(I-\Pi_a)\widetilde G_a(I-\Pi_a^{\mathrm c}).
	\]
	右端先投到值域，\(\widetilde G_a\) 在此给一个解，最后减去核分量。因此 \(T_aG_af=(I-\Pi_a^{\mathrm c})f\)。对 \(u\in X\)，\(G_aT_au\) 与 \(u\) 相差核元素，又与核正交，所以等于 \(u-\Pi_au\)。另外两个恒等式同样由构造成立。若另一算子也满足这四式，对每个 \(f\) 所得两个输出之差同时在核与核正交补中，因而为零。局部构造遂在重叠处一致，给出整个 \(U\) 上的 \(C^r\) 族。

	由 \(u=G_aT_au+\Pi_au\) 得
	\[
		\|u\|_X\le C_a\bigl(\|T_au\|_{H_1}+\|u\|_{H_0}\bigr).
	\]
	若 \(u_j\to u\) 于 \(H_0\) 且 \(T_au_j\to f\) 于 \(H_1\)，此界使 \(u_j\) 在 \(X\) 中为 Cauchy 列；极限只能是 \(u\)，且 \(T_au=f\)。所以 \(\mathsf T_a\) 闭。最后，\(v\perp\operatorname{ran}T_a\) 等价于
	\(\langle T_au,v\rangle=0=\langle u,0\rangle\) 对所有 \(u\in X\) 成立；这正是 \(v\in\operatorname{dom}\mathsf T_a^*\)、\(\mathsf T_a^*v=0\) 的定义。
\end{proof}

在边值问题中，\(\ker\mathsf T_a^*\) 由伴随微分表达式和伴随边界条件共同决定。上面的证明不把它误写成只满足内部伴随方程的解空间。另一方面，若把目标改成 \(H^k\) 并使用该空间的内积，得到的是另一套正交条件。几何上的 \(L^2\) 规范解应先在 \(X=H^2_B\)、\(H_0=H_1=L^2\) 中构造，再由边界正则性提升；下标 \(B\) 表示已经固定的齐次边界条件。

\subsection{参数椭圆估计和连续性}
\label{g3:subsec:1404:02}
一个算子拥有椭圆估计，还不能直接说明整族算子共享同一常数。固定定义域后，吸收一个小的算子扰动就给出局部统一常数；紧参数集再将局部结论拼成整体结论。

\begin{proposition}\label{g3:prop:uniform-parameter-elliptic}
	设 \(X,Y,Z\) 为 Banach 空间，\(X\hookrightarrow Z\) 连续，\(U\subset\mathbb R^s\) 开，\(a\mapsto T_a\in\mathcal L(X,Y)\) 连续。若在 \(a_0\in U\) 有
	\[
		\|u\|_X\le C_0\bigl(\|T_{a_0}u\|_Y+\|u\|_Z\bigr)
		\quad(u\in X),
	\]
	则在 \(a_0\) 的某邻域内，同一估计对 \(T_a\) 成立，常数可取 \(2C_0\)。若每个参数都满足这样的估计，则任意紧集 \(K\subset U\) 上有统一常数。
\end{proposition}
\begin{proof}
	令 \(\|T_a-T_{a_0}\|_{X\to Y}\le(2C_0)^{-1}\)。三角不等式给出
	\[
		\|u\|_X\le C_0\bigl(\|T_au\|_Y+\|u\|_Z\bigr)
		+\tfrac12\|u\|_X.
	\]
	吸收最后一项即得结论。对紧集取这些邻域的有限子覆盖，再取有限个常数的最大值。
\end{proof}

例如，在固定光滑有界域 \(\Omega\subset\mathbb R^m\) 上令
\[
	P_au=-\partial_i(a_a^{ij}\partial_ju)+c_au,
	\qquad X_k=H^{k+2}(\Omega)\cap H^1_0(\Omega),\quad Y_k=H^k(\Omega).
\]
这里 \(k\ge0\) 为整数，每个参数的系数均光滑到边界，矩阵 \((a_a^{ij})\) 实对称且一致正定，\(c_a\) 为实值光滑函数；本段不要求 \(c_a\ge0\)。这里与以后均对重复的 \(i,j\) 求和。若主部系数随参数在 \(C^{k+1}(\overline\Omega)\) 中连续，零阶系数在 \(C^k(\overline\Omega)\) 中连续，则乘积求导给出
\begin{equation}\label{g3:eq:parameter-coefficient-norm}
	\|(P_a-P_b)u\|_{H^k}
	\le C\left(\sum_{i,j}\|a_a^{ij}-a_b^{ij}\|_{C^{k+1}}
	+\|c_a-c_b\|_{C^k}\right)\|u\|_{H^{k+2}}.
\end{equation}
具体说，\(|\alpha|\le k\) 时，\(\partial^\alpha\partial_i(a^{ij}\partial_ju)\) 的每项都是系数的至多 \(k+1\) 阶导数乘以 \(u\) 的至多 \(k+2\) 阶导数；用系数的上确界和后者的 \(L^2\) 范数即可。于是\cref{g3:thm:boundary-higher-regularity}的零 Dirichlet 边界估计在紧参数集上给出
\begin{equation}\label{g3:eq:uniform-dirichlet-estimate}
	\|u\|_{H^{k+2}}\le C_K
	\bigl(\|P_au\|_{H^k}+\|u\|_2\bigr),\qquad a\in K,\quad u\in X_k.
\end{equation}
若系数按上述范数为 \(C^r\)，式\eqref{g3:eq:parameter-coefficient-norm}也说明 \(a\mapsto P_a\) 按 \(X_k\to Y_k\) 的范数为 \(C^r\)。

\begin{corollary}\label{g3:cor:uniform-parameter-green}
	设 \(H_0,H_1\) 为 Hilbert 空间，Banach 空间 \(X\) 连续稠密嵌入 \(H_0\)，\(U\subset\mathbb R^s\) 开，\(a\mapsto T_a\in\mathcal L(X,H_1)\) 为 \(C^r\) 的 Fredholm 族，\(0\le r\le\infty\)，核维数局部恒定。令 \(\Pi_a,\Pi_a^{\mathrm c}\) 分别为核及值域正交补的正交投影，\(G_a\) 为满足 \(T_aG_a=I-\Pi_a^{\mathrm c}\)、\(G_aT_a=I-\Pi_a\)、\(\Pi_aG_a=G_a\Pi_a^{\mathrm c}=0\) 的 Green 算子。设 \(K\subset U\) 紧。则有 \(C_K<\infty\)，使
	\[
		\|G_af\|_X\le C_K\|f\|_{H_1},\qquad
		\|u\|_X\le C_K\|T_au\|_{H_1}\quad(\Pi_au=0),\qquad a\in K.
	\]
	若 \(f_a\in H_1\) 为 \(C^r\) 族且 \(\Pi_a^{\mathrm c}f_a=0\)，则规范解 \(u_a=G_af_a\) 在 \(X\) 中为 \(C^r\)。当 \(r\ge1\)、参数为 \(t\) 时，其导数满足
	\begin{equation}\label{g3:eq:normalized-solution-derivative}
		\dot u_t=G_t(\dot f_t-\dot T_tu_t)-\dot\Pi_tu_t.
	\end{equation}
\end{corollary}
\begin{proof}
	连续映射 \(a\mapsto\|G_a\|\) 在紧集上有界，这给出第一式；对核正交的 \(u\) 用 \(u=G_aT_au\) 即得第二式。算子与向量的配对是连续双线性映射，所以 \(u_a\) 为 \(C^r\)。微分 \(T_tu_t=f_t\) 得
	\(T_t\dot u_t=\dot f_t-\dot T_tu_t\)，微分 \(\Pi_tu_t=0\) 得
	\(\Pi_t\dot u_t=-\dot\Pi_tu_t\)。将两式代入
	\(\dot u_t=G_tT_t\dot u_t+\Pi_t\dot u_t\)，便得导数公式。
\end{proof}

导数公式中的最后一项反映规范条件本身正在改变。核固定时投影固定，这一项为零；核变化时一般需要保留。统一椭圆估计控制的是含低阶范数的左逆；去掉低阶项所需的核正交估计，还要控制零附近的谱隙。

\subsection{核维数跳跃及谱流动机}
\label{g3:subsec:1404:03}
前面分块证明在不假设恒核维数时仍给出
\[
	\dim\ker T_a=\dim\ker S_a\le\dim\ker T_{a_0}.
\]
这称为核维数的上半连续性：某个参数处多出的齐次解，离开该参数便可能消失。指数保持不变，因为核和余核的变化都由同一个有限维映射 \(S_a\) 控制。

\begin{example}\label{g3:exm:dirichlet-resonance-parameter}
	取实值函数空间 \(X=H^2(0,\pi)\cap H^1_0(0,\pi)\)，并令
	\[
		P_tu=-u''-tu,\qquad |t-1|<\tfrac12.
	\]
	齐次方程的解为 \(A\sin(\sqrt t\,x)+B\cos(\sqrt t\,x)\)。两个端点条件先给 \(B=0\)，再给 \(A\sin(\sqrt t\,\pi)=0\)。在指定参数区间中，核仅在 \(t=1\) 为 \(\mathbb R\sin x\)，其余参数核为零。

	每个 \(P_t:X\to L^2\) 都是 Fredholm 的。对此可在 \(t=1\) 直接取归一化正弦基 \(e_j=\sqrt{2/\pi}\sin(jx)\)：它由圆周 Fourier 基的奇延拓得到，且 \(P_te_j=(j^2-t)e_j\)。在去掉可能的 \(j=1\) 模式后，除以 \(j^2-t\) 给出有界的 \(L^2\to H^2\cap H^1_0\) 逆，因为 \((1+j^2)/|j^2-t|\) 统一有界。级数按 \(H^2\) 收敛，可逐项求导到二阶，并保留零端点迹；故它确实解出全部与核正交的右端。

	对 \(t\ne1\)，右端 \(e_1\) 的解为 \(u_t=(1-t)^{-1}e_1\)。因此逆的 \(L^2\to L^2\) 范数至少为 \(|1-t|^{-1}\)。核投影在 \(t\ne1\) 为零，在 \(t=1\) 却为秩一投影，算子范数上不连续。主符号一直是 \(\xi^2\)，而式\eqref{g3:eq:uniform-dirichlet-estimate}的常数仍可统一选取。发散来自趋近零的特征值，不是椭圆性的丧失。
\end{example}

核维数只记录“是否等于零”，没有记录特征值穿过零的方向。有限维 Hermitian 矩阵给出\term[puliu]{谱流}[spectral flow]最直接的模型。对连续道路 \(A:[\alpha,\beta]\to\operatorname{Herm}(N)\)，假设两个端点可逆，令 \(n_-(A)\) 为负特征值个数，按重数计，定义
\begin{equation}\label{g3:eq:finite-spectral-flow}
	\operatorname{sf}(A)=n_-(A(\alpha))-n_-(A(\beta)).
\end{equation}
因此从负侧穿向正侧计 \(+1\)。这个有限维定义不要求选择一组连续的特征向量。

\begin{proposition}\label{g3:prop:spectral-crossing-finite}
	设 \(A:[\alpha,\beta]\to\operatorname{Herm}(N)\) 为 \(C^1\) 道路，端点可逆。若每个满足 \(K_t=\ker A(t)\ne0\) 的内点 \(t\) 上，Hermitian 形式
	\[
		\Gamma_t(v,w)=\langle\dot A(t)v,w\rangle,\qquad v,w\in K_t,
	\]
	均非退化，则这些交叉时刻有限，且
	\[
		\operatorname{sf}(A)=\sum_{t:\,K_t\ne0}\operatorname{sign}\Gamma_t.
	\]
	这里 \(\operatorname{sf}(A)=n_-(A(\alpha))-n_-(A(\beta))\)，\(n_-\) 为负特征值按重数计数；\(\operatorname{sign}\Gamma_t\) 为正惯性指数减负惯性指数。
\end{proposition}
\begin{proof}
	固定交叉时刻 \(t_0\)，按 \(\mathbb C^N=K_{t_0}^\perp\oplus K_{t_0}\) 写
	\[
		A(t)=\begin{pmatrix}B(t)&C(t)\\ C(t)^*&D(t)\end{pmatrix}.
	\]
	\(B(t_0)\) 可逆，而 \(C(t_0)=D(t_0)=0\)。令
	\(S(t)=D(t)-C(t)^*B(t)^{-1}C(t)\)。对二次型完成平方得到
	\[
		\langle A(t)(x,y),(x,y)\rangle
		=\langle B(t)(x+B(t)^{-1}C(t)y),x+B(t)^{-1}C(t)y\rangle
		+\langle S(t)y,y\rangle.
	\]
	可逆坐标变换不改变最大负定子空间的维数，所以
	\(n_-(A(t))=n_-(B(t))+n_-(S(t))\)。这是二次型惯性定理的直接用法。又有
	\[
		S(t)=(t-t_0)\Gamma_{t_0}+o(|t-t_0|),
	\]
	其中将 \(\Gamma_{t_0}\) 用所选内积表示为 Hermitian 算子。

	可逆 Hermitian 矩阵的惯性在充分小的扰动下不变：在其正、负特征子空间上，二次型分别有一个严格的正下界、负上界，小扰动保留这两个界；这两个子空间的维数相加已是总维数。应用此事实，\(B(t)\) 的负惯性局部不变，而当 \(t\) 位于 \(t_0\) 两侧且充分接近时，\(S(t)\) 的惯性分别等于 \(-\Gamma_{t_0}\)、\(\Gamma_{t_0}\) 的惯性。因此该处负特征值数的减少量为
	\[
		n_-(-\Gamma_{t_0})-n_-(\Gamma_{t_0})
		=n_+(\Gamma_{t_0})-n_-(\Gamma_{t_0}).
	\]
	同时 \(S(t)\) 在删去 \(t_0\) 的这个邻域内可逆，故交叉孤立。交叉集合是紧区间中的闭集，且不包含端点；若它无限，就有一个积点，积点本身也是交叉，与孤立性矛盾。所以交叉有限。其余区间上负惯性不变，将各处的减少量相加即得公式。
\end{proof}

例如，\(A(t)=\operatorname{diag}(t,-2t,1+t^2)\)、\(-1\le t\le1\)，在零点有二维核，交叉形式为 \(\operatorname{diag}(1,-2)\)，签名为零。两个特征值向相反方向穿越，谱流为零。若改成 \(\operatorname{diag}(t,t,1+t^2)\)，谱流为二。它们的零点核维数相同，穿越信息却不同。

\ProseParagraphBreak

第十三章的圆周族\cref{g3:exm:kernel-jump-circle}进一步说明无限维情形的动机。对 \(P_t=-\mathrm i\partial_x+t\)、\(-1/2\le t\le1/2\)，固定域为 \(H^1(S^1;\mathbb C)\)，Fourier 模式 \(e^{\mathrm inx}\) 的特征值为 \(n+t\)。其中只有常数模式穿过零，方向为正，净穿越数为一；然而每个端点都有无穷多个负特征值，不能将式\eqref{g3:eq:finite-spectral-flow}的两个有限整数直接换成两个无穷数。

一般谱流理论的一个常用框架是：在固定复 Hilbert 空间 \(H\) 上取闭自伴 Fredholm 算子道路，端点可逆，并要求其闭图 \(\{(u,A_tu):u\in\operatorname{dom}A_t\}\subset H\oplus H\) 上的正交投影按算子范数连续。这就是图的 gap 连续性；它允许定义域变化。在这一框架中，一般理论用零附近的有限维谱空间拼接净穿越数。这里已经证明的是有限维公式，并算出了一个具有显式 Fourier 分解的无限维模型；一般道路的谱投影拼接、选择无关性和同伦不变性是进一步学习谱流时需要建立的内容。

\subsection{参数化求解与几何族}
\label{g3:subsec:1404:04}
现在将参数和非齐次边界数据一起放入求解问题。给出一个内部延拓既能准确表达边界值，也能使空间保持固定。

\begin{theorem}\label{g3:thm:dirichlet-parameter-solution}
	设 \(\Omega\subset\mathbb R^m\) 为非空有界光滑域，\(U\subset\mathbb R^s\) 开，\(k\ge0\) 为整数，\(0\le r\le\infty\)。对每个 \(a\in U\)，令
	\[
		P_au=-\partial_i(a_a^{ij}\partial_ju)+c_au,
	\]
	其中系数光滑到边界，\((a_a^{ij})\) 实对称且在 \(\overline\Omega\) 上一致正定，\(c_a\ge0\)；主部系数随参数在 \(C^{k+1}(\overline\Omega)\) 中为 \(C^r\)，\(c_a\) 在 \(C^k(\overline\Omega)\) 中为 \(C^r\)。设 \(f_a\in H^k(\Omega)\)、\(w_a\in H^{k+2}(\Omega)\) 为 \(C^r\) 族。则有唯一的 \(u_a\in H^{k+2}(\Omega)\)，满足
	\[
		P_au_a=f_a,\qquad \gamma u_a=\gamma w_a,
	\]
	且 \(a\mapsto u_a\) 在 \(H^{k+2}\) 中为 \(C^r\)。对每个紧集 \(K\subset U\)，有统一估计
	\[
		\|u_a\|_{H^{k+2}}\le C_K
		\bigl(\|f_a\|_{H^k}+\|w_a\|_{H^{k+2}}\bigr),\qquad a\in K.
	\]
	解只依赖于 \(\gamma w_a\)，不依赖所选延拓。若 \(r\ge1\)，则一参数导数由
	\begin{equation}\label{g3:eq:dirichlet-parameter-derivative}
		P_t\dot u_t=\dot f_t-\dot P_tu_t,
		\qquad \gamma\dot u_t=\gamma\dot w_t
	\end{equation}
	唯一确定。
\end{theorem}
\begin{proof}
	固定 \(a\)。在 \(H^1_0\) 中，弱形式
	\[
		\mathfrak a_a(v,v)=\int_\Omega
		\left(a_a^{ij}\partial_jv\,\overline{\partial_iv}+c_a|v|^2\right)
	\]
	由一致正定性和\cref{g3:lem:boundary-poincare}的 Dirichlet Poincar\'e 不等式控制 \(\|v\|_{H^1}^2\)。因此\cref{g3:thm:weak-general-dirichlet}给出任意 \(L^2\) 右端的唯一弱解；\cref{g3:thm:boundary-higher-regularity}将 \(H^k\) 右端的解提升为 \(H^{k+2}\)。齐次解以自身测试即为零，所以
	\[
		T_a=P_a|_{X_k}:X_k=H^{k+2}\cap H^1_0\longrightarrow H^k
	\]
	是有界双射；开映射定理保证其逆有界。

	式\eqref{g3:eq:parameter-coefficient-norm}证明 \(a\mapsto T_a\) 为 \(C^r\)，故局部 Neumann 级数及逆映射的光滑性证明 \(a\mapsto T_a^{-1}\) 为 \(C^r\)。令
	\begin{equation}\label{g3:eq:parameter-boundary-lifting}
		u_a=w_a+T_a^{-1}(f_a-P_aw_a).
	\end{equation}
	这给出所述方程、边界条件和 \(C^r\) 依赖。若另一个延拓为 \(w_a+h_a\)，则 \(h_a\in X_k\)，公式的变化为
	\(h_a-T_a^{-1}P_ah_a=0\)，所以解与选择无关。

	紧参数集上，系数范数、\(\|T_a^{-1}\|_{H^k\to X_k}\) 均有界，代入式\eqref{g3:eq:parameter-boundary-lifting}即得估计。最后在固定空间中微分方程，连续迹算子与参数微分可交换，得到式\eqref{g3:eq:dirichlet-parameter-derivative}。同样的齐次唯一性使它唯一确定 \(\dot u_t\)。
\end{proof}

例如，\(w_a(x)=a_1+a_2x_1\) 给出一个明确的光滑边界数据族。一般地，若直接给定 \(C^r\) 族 \(g_a\in H^{k+3/2}(\partial\Omega)\)，则\cref{g3:thm:half-integer-boundary-trace}给出固定的有界线性右逆 \(E_D:H^{k+3/2}(\partial\Omega)\to H^{k+2}(\Omega)\)。取 \(w_a=E_Dg_a\)，上一结果即适用，并将估计中的 \(\|w_a\|_{H^{k+2}}\) 换成 \(C\|g_a\|_{H^{k+3/2}}\)。固定右逆的线性连续性使延拓自动保持参数的 \(C^r\) 正则性。

\ProseParagraphBreak

几何族还需要把不同纤维上的函数识别。设 \(M\) 为固定紧光滑流形，允许有光滑边界，\(\Phi_a:M\to M_a\) 是一族将边界映到边界的光滑微分同胚。给 \(M_a\) 配度量 \(g_a\)，拉回后的密度相对固定正密度 \(\mu_0\) 写为
\[
	\Phi_a^*\mu_{g_a}=\rho_a\mu_0,\qquad \rho_a>0.
\]
若全部拉回系数与 \(\rho_a\) 对参数光滑到所需阶，则
\begin{equation}\label{g3:eq:geometric-unitary-identification}
	U_a:L^2(M_a,\mu_{g_a})\longrightarrow L^2(M,\mu_0),
	\qquad U_av=\rho_a^{1/2}(v\circ\Phi_a)
\end{equation}
是酉同构。事实上换元公式给出
\(\|U_av\|_{L^2(\mu_0)}^2=\int_M|v\circ\Phi_a|^2\rho_a\mu_0=\|v\|_{L^2(\mu_{g_a})}^2\)，逆由 \(\Phi_a^{-1}\) 和 \(\rho_a^{-1/2}\) 明确给出。单纯拉回已经统一底流形，乘上密度平方根又统一了内积。

\ProseParagraphBreak

以 \(P_a=\Delta_{g_a}+1\) 的 Dirichlet 实现为例。拉回坐标中的系数为
\[
	\Delta_{\Phi_a^*g_a}v
	=-|\widetilde g_a|^{-1/2}\partial_i
	\left(|\widetilde g_a|^{1/2}\widetilde g_a^{ij}\partial_jv\right),
	\qquad \widetilde g_a=\Phi_a^*g_a.
\]
乘以正光滑函数保持零迹，因此
\[
	U_a\bigl(H^2(M_a)\cap H^1_0(M_a)\bigr)
	=H^2(M)\cap H^1_0(M)=X.
\]
变换后的 \(\widehat P_a=U_aP_aU_a^{-1}\) 具有固定域 \(X\) 与固定值域 \(L^2(M,\mu_0)\)。坐标公式、乘积求导和有限坐标覆盖证明 \(a\mapsto\widehat P_a\in\mathcal L(X,L^2)\) 光滑；主部仍正定。原方程的弱形式为
\(\int_{M_a}(|\mathrm du|_{g_a}^2+|u|^2)\mu_{g_a}\)，故弱存在和边界正则性给出双射。于是 \(\widehat P_a^{-1}\) 光滑，拉回并调整密度后的解光滑依赖参数。非齐次数据先选择同样已识别的延拓 \(\widehat w_a\)，再用式\eqref{g3:eq:parameter-boundary-lifting}；所有导数均在固定 \(M\) 上计算。

\ProseParagraphBreak

Neumann 与 Hodge 边界条件还包含法向或法向导数，识别底流形后需继续检查定义域。例如在固定区域、仅改变密度的简单情形，写 \(u=\rho_a^{-1/2}v\)，则
\[
	\partial_\nu u=0
	\quad\Longleftrightarrow\quad
	\partial_\nu v-\tfrac12(\partial_\nu\log\rho_a)v=0.
\]
调整密度已经将零 Neumann 条件变成随参数变化的 Robin 条件；若度量也变，法向本身还会改变。要对这样的几何族应用本节定理，须先给出有界同构 \(J_a:X_0\to X_a\) 来识别定义域，并证明搬回的算子 \(U_aP_aJ_a:X_0\to H_1\) 具有所需算子范数正则性。若还要搬回核的正交投影，也须核对这些识别与底层内积的关系。Hodge 的绝对、相对条件同样要检查法向分量与相应导数。固定的边界条件名称本身不能代替这个定义域识别。
% 纲要条目排为编号小节标题，并进入目录与书签；说明段仍保留于纲要注释中。


\begin{chaptersummary}
	Green 公式把内部方程与边界通量联系起来。Dirichlet 条件通过仿射迹类限制未知量及测试，Neumann 条件通过边界对偶进入右端；常数核对应逐连通分支的总通量兼容条件。混合条件仍可给出强制的弱问题，而角点与条件交界的显式解说明，弱存在性和整体二阶正则性需要分别论证。

	\ProseParagraphBreak

	光滑边界上的正则性先由切向差商建立，再用方程恢复法向导数。半整数阶迹与有界提升把非齐次边界数据转化为齐次问题，高阶估计记录了内部右端与两类边界数据不同的导数阶数。对于系统，衰减模型中的补条件解释了边界算子应当控制哪些方向；绝对、相对形式条件把这一结构与能量估计结合起来，得到有限维调和核、闭值域及相应的规范解。

	\ProseParagraphBreak

	参数依赖需要先把变化的对象识别到固定空间。Fredholm 块分解在恒定核维数下给出连续或光滑的核投影与广义逆，统一的椭圆估计再控制高阶范数。核跳跃处，有限维剩余方程记录共振与逆的发散；特征值穿零的具体模型则引出谱流所计数的现象。几何族中的拉回、密度和边界域都是这种识别的一部分。
\end{chaptersummary}

\BookExercises

\begin{exercise}\label{g3:ex:interval-variable-flux}
	设 \(a\in C^\infty([0,L];\mathbb R)\)、\(a>0\)，\(f\in L^2(0,L)\)，考虑 \(-(au')'=f\)。
	\begin{enumerate}
		\item 给定 \(u(0)=\alpha\)、\(u(L)=\beta\)，写出解的积分公式。
		\item 若改为指定外通量 \(-a(0)u'(0)=g_0\)、\(a(L)u'(L)=g_L\)，求兼容条件及唯一零均值解。
	\end{enumerate}
\end{exercise}
\begin{solution}
	（1）令 \(F(x)=\int_0^xf(t)\,\mathrm dt\)。积分方程得到 \(a(x)u'(x)=J-F(x)\)，所以
	\[
		u(x)=\alpha+\int_0^x\frac{J-F(s)}{a(s)}\,\mathrm ds,\qquad
		J=\frac{\beta-\alpha+\int_0^LF(s)/a(s)\,\mathrm ds}
		{\int_0^L1/a(s)\,\mathrm ds}.
	\]
	分母为正，因此没有额外兼容条件。所得函数属于 \(H^2\)，分部积分验证弱方程和端点值。

	（2）左端通量使 \(J=-g_0\)，右端要求 \(-g_0-F(L)=g_L\)，即
	\(\int_0^Lf+g_0+g_L=0\)。条件成立时全部解为
	\[
		u(x)=C-\int_0^x\frac{g_0+F(s)}{a(s)}\,\mathrm ds.
	\]
	交换积分次序，零均值条件给出
	\[
		C=\frac1L\int_0^L\frac{(L-s)(g_0+F(s))}{a(s)}\,\mathrm ds.
	\]
	两种问题的齐次差均满足 \(a u'=\text{常数}\)。零端点值使该常数及函数同时为零；零通量只使导数为零，仍需均值确定函数常数。
\end{solution}

\begin{exercise}\label{g3:ex:ball-neumann-flux-sign}
	在 Euclid 单位球 \(B^n\)、\(n\ge1\)，求 \(\Delta u=-2n\)、\(\partial_\nu u=2\) 的零均值解。用体积与边界面积直接验证兼容条件，并说明若将外通量改为零会怎样。这里 \(\Delta=-\sum_i\partial_i^2\)。
\end{exercise}
\begin{solution}
	\(u_0(x)=|x|^2\) 满足 \(\Delta u_0=-2n\)，球面上 \(\partial_\nu u_0=2\)。记单位球体积为 \(\omega_n\)，边界测度为 \(n\omega_n\)，则
	\[
		\int_{B^n}(-2n)+\int_{\partial B^n}2
		=-2n\omega_n+2n\omega_n=0.
	\]
	径向积分给出 \(\int_{B^n}|x|^2=n\omega_n/(n+2)\)，所以唯一零均值解为
	\[
		u(x)=|x|^2-\frac n{n+2}.
	\]
	\(n=1\) 时边界是两个端点，计数测度给出相同公式。若通量改成零，兼容条件的左侧为 \(-2n\omega_n\ne0\)，因而连弱解也不存在；这与内部多项式原函数是否存在无关。
\end{solution}

\begin{exercise}\label{g3:ex:disk-dirichlet-neumann-map}
	在单位圆盘上，给定边界数据 \(\psi(\theta)=\sum_{k\in\mathbb Z}c_ke^{\mathrm ik\theta}\)，满足 \(\sum_k(1+|k|)|c_k|^2<\infty\)。用圆盘内的调和函数计算从 Dirichlet 数据到外法向通量的映射，并证明
	\[
		\int_{B^2}|\mathrm du|^2=2\pi\sum_k|k||c_k|^2.
	\]
	说明其核及通量的平均值；同时解释为什么所得通量总能作用于 \(H^{1/2}(S^1)\)，却未必属于 \(L^2(S^1)\)。
\end{exercise}
\begin{solution}
	先说明圆周半阶空间的 Fourier 表达。对周期函数，半阶差分积分可取 \(\int_{-\pi}^{\pi}\int_0^{2\pi}|\psi(\theta+h)-\psi(\theta)|^2|h|^{-2}\,\mathrm d\theta\mathrm dh\)。它与局部图中的半阶范数等价：小的 \(h\) 对应同图中两点的距离，大的 \(h\) 部分由 \(L^2\) 范数控制。Parseval 将该积分写成各频率平方乘以 \(2\pi\int_{-\pi}^{\pi}|e^{\mathrm ikh}-1|^2/h^2\,\mathrm dh\)。当 \(k\ne0\) 时令 \(s=|k|h\)，最后的积分与 \(|k|\) 一致可比，因为被积函数在零附近有界、在无穷远为 \(O(s^{-2})\)，且在 \([-\pi,\pi]\) 上积分严格为正。因此题设平方和正是 \(H^{1/2}(S^1)\) 的一个等价范数。

	对有限 Fourier 和取
	\[
		u(r,\theta)=\sum_kc_kr^{|k|}e^{\mathrm ik\theta}.
	\]
	每项在原点也是光滑调和多项式。不同角频率在圆周上正交；对 \(k\ne0\)，径向和角向梯度项各贡献 \(\pi|k||c_k|^2\)，故得到题给能量等式。另外
	\[
		\|u\|_{L^2(B^2)}^2=\pi\sum_k\frac{|c_k|^2}{|k|+1}.
	\]
	这两个公式使部分和在 \(H^1(B^2)\) 中收敛；其迹在边界 \(L^2\) 中趋于 \(\psi\)。弱调和性由测试配对传到极限，所以这是唯一的 Dirichlet 弱解。其弱通量的 Fourier 系数为 \(|k|c_k\)，即
	\[
		\Lambda\psi=\sum_k|k|c_ke^{\mathrm ik\theta}.
	\]
	若 \(\varphi=\sum_kd_ke^{\mathrm ik\theta}\in H^{1/2}\)，则
	\[
		|\langle\Lambda\psi,\varphi\rangle|
		=\left|2\pi\sum_k|k|c_k\overline{d_k}\right|
		\le2\pi\left(\sum_k|k||c_k|^2\right)^{1/2}
		\left(\sum_k|k||d_k|^2\right)^{1/2}.
	\]
	以 \(H^{-1/2}\) 表示 \(H^{1/2}\) 的连续反线性对偶，这个估计说明 \(\Lambda:H^{1/2}\to H^{-1/2}\) 有界。零频率消失，所以通量平均值为零；核恰为常数。取 \(c_k=k^{-3/2}\) 于 \(k\ge1\)、其余为零，则边界数据满足题设，而通量的平方和为 \(\sum_{k\ge1}k^{-1}=\infty\)，故通量不属于 \(L^2\)。
\end{solution}

\begin{exercise}\label{g3:ex:positive-boundary-energy}
	设 \(M\) 为连通紧光滑带边界黎曼流形，\(\partial M\ne\varnothing\)，\(\beta>0\) 为常数。证明
	\[
		Q_\beta(u,v)=\int_M\langle\mathrm du,\mathrm dv\rangle\,\mu_g
		+\beta\int_{\partial M}\gamma u\,\overline{\gamma v}\,\mu_{\partial}
	\]
	在 \(H^1(M)\) 上强制，并据此求解弱边值问题
	\(\Delta u=f\)、\(\partial_\nu u+\beta\gamma u=g\)，其中 \(f\in L^2(M)\)、\(g\in L^2(\partial M)\)。在区间 \([0,1]\) 上找出 \(\beta<0\) 时使齐次问题出现非零解的全部常数 \(\beta\)。
\end{exercise}
\begin{solution}
	若不存在强制常数，可取 \(\|u_j\|_{H^1}=1\)，而 \(\|\mathrm du_j\|_2\to0\)、\(\|\gamma u_j\|_2\to0\)。紧嵌入与弱紧性给出 \(u_j\to u\) 于 \(L^2\)、弱收敛于 \(H^1\)。因梯度范数趋零，有 \(\|u\|_2=1\)，且 \(u\) 为常数。迹的弱连续性又给出 \(\gamma u=0\)，边界非空使该常数为零，矛盾。故 \(Q_\beta\) 强制。

	迹估计使右端 \((f,v)_2+\int_{\partial M}g\overline{\gamma v}\) 有界。Lax--Milgram 给出唯一 \(u\in H^1\)。内部测试给出 \(\Delta u=f\)；弱共法向定义给出 \(\gamma_Iu+\beta\gamma u=g\)，这就是所指定边界条件的弱含义。

	区间上齐次方程使 \(u(x)=a x+b\)。两端的外法向相反，故条件为
	\[
		-a+\beta b=0,\qquad a+\beta(a+b)=0.
	\]
	第一式给 \(a=\beta b\)，第二式化为 \(\beta(\beta+2)b=0\)。非零解且 \(\beta<0\) 时只能 \(\beta=-2\)，解为 \(u=b(1-2x)\)。负边界项可以抵消内部梯度能量，从而产生非平凡核。
\end{solution}

\begin{exercise}\label{g3:ex:conormal-not-normal}
	在平面半域 \(\{(x,t):t>0\}\) 上取
	\[
	A=\begin{pmatrix}2&1\\1&2\end{pmatrix},\qquad
	P=-\operatorname{div}(A\nabla).
	\]
	\begin{enumerate}
		\item 对 \(u(x,t)=x-\tfrac12t\)，计算 \(Pu\)、\(\partial_\nu u\) 与 \(B_\nu u\)，其中平面边界的外法向量为 \(\nu=-e_t\)。
		\item 设 \(g\in C_c^\infty(\mathbb R)\)。构造在 \(t=0\) 附近满足
		\(\gamma w=0\)、\(B_\nu w=g\) 的光滑函数，并说明为何一般提升时需除以 \(\nu\cdot A\nu\)。
	\end{enumerate}
\end{exercise}
\begin{solution}
	线性函数的二阶导数为零，故 \(Pu=0\)。又
	\[
	\nabla u=(1,-\tfrac12),\qquad
	A\nabla u=(\tfrac32,0),
	\]
	所以 \(\partial_\nu u=-u_t=\tfrac12\)，而 \(B_\nu u=0\)。因此齐次共法向条件并不等于齐次普通法向条件。

	取 \(\eta\in C_c^\infty([0,\infty))\)，使其在零点附近为一，令
	\[
	w(x,t)=-\tfrac12t\,g(x)\eta(t).
	\]
	在 \(t=0\)，\(w=0\)、\(w_x=0\)、\(w_t=-g/2\)，从而
	\(B_\nu w=-(w_x+2w_t)=g\)。一般情形下，零迹使所有切向导数迹为零，因此边界梯度只有法向分量，且
	\[
	B_\nu w=(\nu\cdot A\nu)\partial_\nu w.
	\]
	这就要求把普通法向导数指定为 \(g/(\nu\cdot A\nu)\)。椭圆性给出分母的正下界。这里显式乘 \(t\) 的公式适用于光滑 \(g\)；要把有限正则性数据提升且只损失规定的半阶，需使用正文的 Fourier 提升。
\end{solution}

\begin{exercise}\label{g3:ex:half-order-trace-scaling}
	在平圆柱 \(C=(\mathbb R/2\pi\mathbb Z)\times(0,1)\) 上取
	\(v_N(x,t)=e^{iNx}e^{-Nt}\)，\(N\ge1\)。固定整数 \(q\ge1\)。
	\begin{enumerate}
		\item 证明 \(\|v_N\|_{H^q(C)}\asymp N^{q-1/2}\)，比较其 \(t=0\) 的函数值与法向导数。
		\item 对周期函数 \(w\in H^q(C)\) 且 \(\gamma_0w=e^{iNx}\)，证明
		\(\|w\|_{H^q(C)}\ge c_qN^{q-1/2}\)。由此解释 \(H^q\) 到边界损失半阶的尺度。
	\end{enumerate}
\end{exercise}
\begin{solution}
	对 \(a+b\le q\)，
	\[
	\|\partial_x^a\partial_t^bv_N\|_{L^2(C)}^2
	=2\pi N^{2(a+b)}\int_0^1e^{-2Nt}\,\mathrm dt
	=\pi N^{2(a+b)-1}(1-e^{-2N}).
	\]
	取总阶数 \(q\) 的项给出下界，对全部有限项求和给出上界。下边界外法向为 \(-\partial_t\)，故两个数据分别为
	\(e^{iNx}\) 和 \(Ne^{iNx}\)。若周期 Sobolev 范数写成
	\(\sum_j(1+j^2)^s|\widehat g_j|^2\)，它们在
	\(H^{q-1/2}\) 与 \(H^{q-3/2}\) 中的范数均与 \(N^{q-1/2}\) 同阶；第二个比较在 \(q\ge2\) 时对应正文的联合迹定理。

	为直接证明任何提升的下界，令
	\[
	W(t)=\frac1{2\pi}\int_0^{2\pi}w(x,t)e^{-iNx}\,\mathrm dx.
	\]
	则 \(W\in H^1(0,1)\)、\(W(0)=1\)。选择固定光滑截断 \(\chi\)，使 \(\chi(0)=1\)、\(\chi=0\) 于 \(t\ge3/4\)。对 \(\chi W\) 使用一维分部积分，并以 \(N\) 配平两个因子，得
	\[
	|W(0)|^2
	\le C\left(N\int_0^1|W|^2\,\mathrm dt
																+N^{-1}\int_0^1|W'|^2\,\mathrm dt\right).
	\]
	乘以 \(N^{2q-1}\)，右侧由
	\(C(N^{2q}\|W\|_2^2+N^{2q-2}\|W'\|_2^2)\) 控制。周期 Fourier 正交性说明这两项不超过常数倍
	\[
	\|\partial_x^qw\|_2^2+
	\|\partial_x^{q-1}\partial_tw\|_2^2.
	\]
	因而 \(\|w\|_{H^q}^2\ge c_qN^{2q-1}\)。光滑逼近使计算对全部 \(H^q\) 函数成立。第一问的 \(v_N\) 达到同阶上界，故这个尺度不能再改善。
\end{solution}

\begin{exercise}\label{g3:ex:rough-neumann-data-harmonic}
	在单位圆盘 \(D\) 上考虑
	\[
	u(r,\theta)=\sum_{j=2}^\infty j^{-7/4}r^j\cos(j\theta).
	\]
	证明该级数定义一个 \(H^1(D)\) 函数，在内部调和，边界迹属于
	\(H^1(\partial D)\)，并且是某个 \(L^2(\partial D)\) Neumann 数据的弱解；但 \(u\notin H^2(D)\)。写出该 Neumann 数据并验证兼容条件。
\end{exercise}
\begin{solution}
	记 \(p_j=r^j\cos(j\theta)\)。不同 \(j\) 的函数、梯度在 \(L^2(D)\) 中分别正交；梯度正交也可由分部积分及
	\(\partial_\nu p_j=j\cos(j\theta)\) 得到。由正文算例，
	\[
	\|p_j\|_2^2=\frac{\pi}{2j+2},\qquad
	\|\nabla p_j\|_2^2=\pi j.
	\]
	因此部分和在 \(H^1(D)\) 中为 Cauchy 列，因为
	\(\sum_{j\ge2}j^{-7/2}(j+(2j+2)^{-1})<\infty\)。
	每个部分和调和，对内部紧支测试函数取 \(H^1\) 极限，得到 \(\Delta u=0\) 的分布等式。内部正则性给出内部光滑性。

	边界级数
	\[
	g_D(\theta)=\sum_{j=2}^\infty j^{-7/4}\cos(j\theta)
	\]
	在 \(H^1(\mathbb R/2\pi\mathbb Z)\) 中收敛，因为
	\(\sum j^2j^{-7/2}=\sum j^{-3/2}<\infty\)。部分和的内部迹收敛到 \(\gamma u\)，故 \(\gamma u=g_D\)。相应外法向数据为
	\[
	g_N(\theta)=\sum_{j=2}^\infty j^{-3/4}\cos(j\theta)\in L^2(\partial D),
	\]
	其平方系数和为 \(\sum j^{-3/2}<\infty\)。对每个部分和 \(u_J\)，分部积分给出
	\[
	\int_D\nabla u_J\cdot\overline{\nabla w}
	=\int_{\partial D}g_{N,J}\,\overline{\gamma w}
	\quad(w\in H^1(D)).
	\]
	左侧按 \(H^1\) 收敛，右侧按 \(L^2\) 与迹连续性收敛，因此 \(u\) 是该 Neumann 数据的弱解。各余弦项的积分为零，\(L^2\) 收敛保留积分，故
	\(\int_D0+\int_{\partial D}g_N=0\)，恰为兼容条件。

	最后证明 \(u\notin H^2(D)\)。若 \(u\in H^2\)，则其二阶导数在较小圆盘上由逐项求导得到。对 \(0<\rho<1\)，不同 \(j\) 的二阶导数矩阵在 \(L^2(D_\rho)\) 中正交，于是
	\[
	\sum_{a,b=1}^2\|\partial_a\partial_bu\|_{L^2(D_\rho)}^2
	=2\pi\sum_{j=2}^\infty j^{-7/2}j^2(j-1)\rho^{2j-2}.
	\]
	当 \(\rho\uparrow1\)，右侧单调趋于无穷，因为各项与 \(j^{-1/2}\) 同阶，矛盾。这里局部逐项求导由几何因子 \(\rho^j\) 保证；不同频率 \(j-2\) 的实部、虚部给出所用正交性。这个例子说明，光滑域上的光滑内部方程配以 \(L^2\) Neumann 数据，仍不足以自动推出 \(H^2\)。
\end{solution}

\begin{exercise}\label{g3:ex:coefficient-regularity-boundary}
	在区间 \((0,1)\) 上令
	\[
	a(t)=1+t^{5/4},\qquad
	u(t)=\int_0^t\frac{\mathrm ds}{a(s)}.
	\]
	证明 \(a\) 一致正且属于 \(C^1([0,1])\)，\(u\) 满足
	\(-\bigl(au'\bigr)'=0\)，边界值与共法向值均为有限的端点数据，但
	\(u\in H^2(0,1)\setminus H^3(0,1)\)。说明这与本节的高阶正则性定理是否冲突。
\end{exercise}
\begin{solution}
	有 \(1\le a(t)\le2\)，且
	\(a'(t)=\tfrac54t^{1/4}\) 连续延伸到零点，所以 \(a\in C^1([0,1])\)。由定义
	\[
	u'=\frac1a,\qquad au'=1,\qquad
	u''=-\frac{a'}{a^2}.
	\]
	前两阶导数都连续有界，故 \(u\in H^2(0,1)\)。Dirichlet 数据为
	\(u(0)=0\)、\(u(1)=\int_0^1a(s)^{-1}\,\mathrm ds\)；
	外法向在左端为 \(-1\)、右端为 \(+1\)，所以共法向数据为
	\(B_\nu u(0)=-1\)、\(B_\nu u(1)=1\)。它们的和为零，与内部源项为零相容。

	在 \(t>0\) 上，
	\[
	u'''(t)
	=-\frac{5}{16}\frac{t^{-3/4}}{(1+t^{5/4})^2}
			+\frac{25}{8}\frac{t^{1/2}}{(1+t^{5/4})^3}.
	\]
	第二项在零点附近有界，第一项的绝对值与 \(t^{-3/4}\) 同阶，其平方不可积。如果 \(u\in H^3(0,1)\)，三阶弱导数在每个内部紧区间上必须等于这个经典导数，从而必须在整个区间平方可积，矛盾。

	高阶定理要求足够高阶的系数正则性；本例的
	\(a''(t)=\tfrac5{16}t^{-3/4}\) 在边界发散，主系数不光滑。因此不存在冲突。二阶结论只用 \(W^{1,\infty}\) 主系数，本例恰好满足并取得这一结论。
\end{solution}

\begin{exercise}\label{g3:ex:complex-oblique-boundary-family}
	设 \(D\subset\mathbb C\) 为单位圆盘，\(a\in\mathbb R\)，对复值函数取
	\[
		P=-\partial_x^2-\partial_y^2,\qquad
		B_a=\partial_\nu+\mathrm i a\partial_\theta.
	\]
	检验哪些 \(a\) 满足椭圆补条件，并确定光滑齐次问题 \(Pu=0,B_au=0\) 的核。解释为什么满足补条件仍不能推出核为零。
\end{exercise}
\begin{solution}
	沿边界取正向弧长坐标 \(y\)，令 \(t\) 指向内部。对衰减模态 \(\mathrm e^{\mathrm i y\xi-|\xi|t}c\)，边界主部给出
	\[
		(|\xi|-a\xi)c.
	\]
	当 \(\xi>0\) 时系数为 \((1-a)\xi\)，当 \(\xi<0\) 时为 \(-(1+a)\xi\)。因此补条件成立当且仅当 \(a\ne1,-1\)。特别地，\(|a|>1\) 仍满足此处没有谱参数的补条件。

	为确定核，把光滑调和函数沿圆周展开。第 \(k\ne0\) 个 Fourier 系数满足径向方程
	\[
		u_k''+r^{-1}u_k'-k^2r^{-2}u_k=0,
	\]
	解为 \(A_kr^{|k|}+B_kr^{-|k|}\)；原点光滑排除第二项。零频率解为 \(A_0+B_0\log r\)，同样排除对数项。因此
	\[
		u(r,\theta)=c_0+\sum_{k\ge1}r^k
		\bigl(c_k\mathrm e^{\mathrm i k\theta}
		+d_k\mathrm e^{-\mathrm i k\theta}\bigr).
	\]
	光滑边界迹保证系数迅速衰减，故可在边界逐项求导。条件变成
	\[
		k(1-a)c_k=0,\qquad k(1+a)d_k=0\quad(k\ge1).
	\]
	若 \(a\ne\pm1\)，核恰为常数；若 \(a=1\)，核为所有在闭圆盘上光滑的全纯函数；若 \(a=-1\)，核为相应的反全纯函数。后两者均无限维。补条件只检验非零切向频率，常数属于它没有排除的全局零频率。
\end{solution}

\begin{exercise}\label{g3:ex:ball-absolute-relative-harmonics}
	设 \(B_R\subset\mathbb R^n\) 为半径 \(R>0\) 的 Euclid 球，\(n\ge2\)。用本节的边界能量恒等式，直接求出所有次数的绝对与相对调和形式空间。这里取复值形式，不调用上同调计算。
\end{exercise}
\begin{solution}
	球面外法向满足 \(S(X)=\nabla_X\nu=R^{-1}X\)，故在切向 \(p\) 次形式上
	\[
		S^{[p]}=\frac pR I,\qquad
		(\operatorname{tr}S)I-S^{[q-1]}=\frac{n-q}{R}I.
	\]
	Euclid 联络曲率为零，因而 \(\mathcal R_q=0\)。由\cref{g3:cor:hodge-boundary-fredholm}，调和核中的形式光滑并满足 \(Q(\omega,\omega)=0\)。绝对情形在 \(1\le q\le n-1\) 时给出
	\[
		0=\|\nabla\omega\|_2^2
		+\frac qR\|\mathrm t\omega\|_{L^2(\partial B_R)}^2.
	\]
	于是 \(\omega\) 平行，切向迹为零；绝对零阶条件又使法向迹为零，所以它在边界为零。平行形式的长度为常数，故 \(\omega=0\)。当 \(q=n\) 时切向迹自动为零，而绝对条件已经使全部边界值为零；同一能量恒等式先给出平行性，仍得零解。\(q=0\) 时能量只要求 \(\mathrm du=0\)，核为常数。因此
	\[
		\mathcal H_{\mathrm A}^0(B_R)=\mathbb C,
		\qquad \mathcal H_{\mathrm A}^q(B_R)=0\quad(1\le q\le n).
	\]

	相对情形在 \(1\le q\le n-1\) 时变成
	\[
		0=\|\nabla\omega\|_2^2
		+\frac{n-q}{R}\|\mathrm n\omega\|_{L^2(\partial B_R)}^2,
	\]
	再结合 \(\mathrm t\omega=0\)，同样得到零解。\(q=0\) 是零 Dirichlet 的常数函数，故也为零。\(q=n\) 时边界系数为零，只得到平行性；所有平行顶次形式是 \(c\,\mathrm dx^1\wedge\cdots\wedge\mathrm dx^n\)，并且确实满足相对条件及 \(\mathrm d\omega=\delta\omega=0\)。所以
	\[
		\mathcal H_{\mathrm R}^q(B_R)=0\quad(0\le q<n),\qquad
		\mathcal H_{\mathrm R}^n(B_R)
		=\mathbb C\,\mathrm dx^1\wedge\cdots\wedge\mathrm dx^n.
	\]
\end{solution}

\begin{exercise}\label{g3:ex:cylinder-hodge-boundary-spectrum}
	设 \(M=[0,L]\times(\mathbb R/2\pi\mathbb Z)\)，带乘积度量 \(\mathrm dt^2+\mathrm d\theta^2\)，\(L>0\)。对一形式 \(\omega=a\,\mathrm dt+b\,\mathrm d\theta\)，把绝对与相对条件写成 \(a,b\) 的端点条件，求两种实现的特征值及一组特征形式，并确定调和核。由此说明式\eqref{g3:eq:hodge-gaffney-estimate}中的 \(L^2\) 项一般不能直接删去。
\end{exercise}
\begin{solution}
	在此平坦乘积中，
	\[
		\mathrm d\omega=(\partial_tb-\partial_\theta a)
		\,\mathrm dt\wedge\mathrm d\theta,
		\qquad
		\delta\omega=-(\partial_ta+\partial_\theta b),
	\]
	而 \(\Delta\) 对两个分量分别作用为 \(-\partial_t^2-\partial_\theta^2\)。端点的外法向分别是 \(-\partial_t\) 与 \(\partial_t\)。绝对零阶条件给出 \(a=0\)；沿边界对该等式求导得到 \(\partial_\theta a=0\)，故一阶条件等价于 \(\partial_tb=0\)。相对零阶条件给出 \(b=0\)，再用 \(\partial_\theta b=0\) 得 \(\partial_ta=0\)。这些等式均在 \(t=0,L\) 上成立。

	绝对实现的特征形式可取
	\[
		\sin\frac{m\pi t}{L}\,\mathrm e^{\mathrm i k\theta}\,\mathrm dt
		\quad(m\ge1,\ k\in\mathbb Z),\qquad
		\cos\frac{m\pi t}{L}\,\mathrm e^{\mathrm i k\theta}\,\mathrm d\theta
		\quad(m\ge0,\ k\in\mathbb Z).
	\]
	相对实现交换正弦与余弦所在的分量：
	\[
		\cos\frac{m\pi t}{L}\,\mathrm e^{\mathrm i k\theta}\,\mathrm dt
		\quad(m\ge0),\qquad
		\sin\frac{m\pi t}{L}\,\mathrm e^{\mathrm i k\theta}\,\mathrm d\theta
		\quad(m\ge1),\qquad k\in\mathbb Z.
	\]
	每个所列形式对应特征值 \((m\pi/L)^2+k^2\)，重数按出现的独立形式计。\(t\) 方向的正弦、余弦系分别由区间上的奇、偶延拓及 Fourier 完备性得到，\(\theta\) 方向使用周期 Fourier 基；它们的乘积因此给出完整的正交基。除以各自的 \(L^2\) 范数便成为正交归一基。

	零特征值要求 \(m=k=0\)，故
	\[
		\mathcal H_{\mathrm A}^1=\mathbb C\,\mathrm d\theta,
		\qquad \mathcal H_{\mathrm R}^1=\mathbb C\,\mathrm dt.
	\]
	这些形式非零而 \(\mathrm d\omega=\delta\omega=0\)。它们的 \(H^1\) 范数含有正的 \(L^2\) 范数，因此若在式\eqref{g3:eq:hodge-gaffney-estimate}中直接删去 \(L^2\) 项，就会得到矛盾。若先限制到相应调和核的正交补，则谱隙配合该估计可恢复不含这一项的强制估计。
\end{solution}

\begin{exercise}\label{g3:ex:robin-negative-eigenvalue}
	给定 \(\kappa>0\)，在 \([0,L]\) 上考虑 \(Pu=-u''\)，端点条件为
	\[
		-u'(0)+\kappa u(0)=0,\qquad
		u'(L)-\kappa u(L)=0.
	\]
	证明该实现自伴、下有界并具有负特征值。写出对应的能量型，说明这与绝对、相对 Hodge 实现的非负性有何关系。
\end{exercise}
\begin{solution}
	两个端点都取自由的 Robin 分量，系数分别为 \(S(0)=\kappa\)、\(S(L)=-\kappa\)。它们是实数，故满足\cref{g3:thm:mixed-bundle-fredholm}的自伴假设；区间是光滑一维带边界流形。因此该实现自伴且下有界，并具有紧预解式。其能量型是
	\[
		a(u,v)=\int_0^L u'\overline{v'}\,\mathrm dt
		+\kappa u(0)\overline{v(0)}
		-\kappa u(L)\overline{v(L)}.
	\]
	直接取 \(u(t)=\mathrm e^{\kappa t}\)，两个端点条件均成立，且
	\[
		Pu=-\kappa^2u.
	\]
	所以 \(-\kappa^2\) 是负特征值。也可直接计算
	\[
		a(u,u)=\frac\kappa2(\mathrm e^{2\kappa L}-1)
		+\kappa-\kappa\mathrm e^{2\kappa L}
		=-\kappa^2\|u\|_{L^2}^2.
	\]
	自伴性保证特征值为实数，并不保证其为非负。Hodge 的绝对与相对实现多了一项具体恒等式：它们的能量型恰为 \(\|\mathrm d\omega\|_2^2+\|\delta\omega\|_2^2\)。非负性来自这个平方和结构，而非一般 Robin 边界条件的名称。
\end{solution}

\begin{exercise}\label{g3:ex:moving-kernel-normalization}
	在标准 Euclidean 空间 \(\mathbb R^2\) 中令
	\[
		T_t=\begin{pmatrix}1&t\\0&0\end{pmatrix},\qquad f=\binom10.
	\]
	求核正交投影 \(\Pi_t\)、余核正交投影 \(\Pi_t^{\mathrm c}\) 与 Green 算子 \(G_t\)。求 \(T_tu_t=f\)、\(u_t\perp\ker T_t\) 的规范解，并验证在 \(t=0\) 处，式\eqref{g3:eq:normalized-solution-derivative}中的移动投影项不能删去。
\end{exercise}
\begin{solution}
	核由 \((-t,1)\) 张成，值域由 \((1,0)\) 张成，所以
	\[
		\Pi_t=\frac1{1+t^2}\begin{pmatrix}t^2&-t\\-t&1\end{pmatrix},
		\qquad\Pi_t^{\mathrm c}=\begin{pmatrix}0&0\\0&1\end{pmatrix}.
	\]
	核正交补由 \((1,t)\) 张成，而 \(T_t(1,t)=(1+t^2,0)\)。先去掉右端的第二坐标再求规范解，得到
	\[
		G_t=\frac1{1+t^2}\begin{pmatrix}1&0\\t&0\end{pmatrix},
		\qquad u_t=\frac{(1,t)}{1+t^2}.
	\]
	直接相乘得 \(T_tG_t=I-\Pi_t^{\mathrm c}\)、\(G_tT_t=I-\Pi_t\)，并且两个规范条件也成立。在零点有 \(u_0=(1,0)\)、\(\dot u_0=(0,1)\)，但
	\[
		\dot T_0u_0=0,\qquad
		\dot\Pi_0=\begin{pmatrix}0&-1\\-1&0\end{pmatrix},\qquad
		-\dot\Pi_0u_0=(0,1).
	\]
	因此导数完全来自核正交补的旋转。若漏掉投影的导数，就会错误地得到 \(\dot u_0=0\)。
\end{solution}

\begin{exercise}\label{g3:ex:principal-coefficient-parameter}
	对 \(|t|<1/2\)，令 \(u_t\in H^2(0,1)\cap H^1_0(0,1)\) 解
	\[
		-\bigl((1+tx)u_t'\bigr)'=1.
	\]
	证明此解存在唯一且光滑依赖 \(t\)。求 \(u_0\) 与 \(w=\partial_tu_t|_{t=0}\)，并证明在 \(H^2(0,1)\) 中有 \(u_t=u_0+tw+O(t^2)\)。
\end{exercise}
\begin{solution}
	主部系数 \(1+tx\ge1/2\)，且参数依赖光滑。\Cref{g3:thm:dirichlet-parameter-solution}给出存在、唯一和 \(H^2\) 中的光滑依赖。直接积分得 \(u_0=x(1-x)/2\)。写 \(P_t=P_0+tB\)，其中 \(P_0u=-u''\)、\(Bu=-(xu')'\)。导数方程为
	\[
		-w''=-Bu_0=(xu_0')'=\tfrac12-2x,\qquad w(0)=w(1)=0.
	\]
	积分两次并用端点条件，得到
	\[
		w(x)=-\frac{x^2}{4}+\frac{x^3}{3}-\frac{x}{12}.
	\]
	令 \(r_t=u_t-u_0-tw\)。因 \(P_0u_0=1\)、\(P_0w+Bu_0=0\)，有
	\[
		P_tr_t=-t^2Bw,\qquad r_t\in H^2\cap H^1_0.
	\]
	在例如 \(|t|\le1/4\) 上，\(P_t^{-1}:L^2\to H^2\cap H^1_0\) 的范数有统一上界。因此
	\(\|r_t\|_{H^2}\le Ct^2\|Bw\|_2\)，得到所需二阶余项。这也直接核对了导数方程中的符号。
\end{solution}

\begin{exercise}\label{g3:ex:moving-interval-dirichlet}
	设 \(L(t)>0\)、\(p(t)\)、\(q(t)\) 为实值 \(C^1\) 函数。求区间 \((0,L(t))\) 上的解
	\[
		-u_t''=1,\qquad u_t(0)=p(t),\quad u_t(L(t))=q(t).
	\]
	用 \(x=L(t)y\) 将方程拉回 \((0,1)\)，分别写出普通拉回 \(v_t(y)=u_t(L(t)y)\) 和酉识别 \(\widehat u_t(y)=L(t)^{1/2}v_t(y)\)。求 \(\dot v_t\) 的方程与端点值，并说明为什么不能将固定 \(y\) 的参数导数当作固定 \(x\) 的参数导数。
\end{exercise}
\begin{solution}
	二次积分给出
	\[
		u_t(x)=p(t)+\frac{q(t)-p(t)}{L(t)}x
		+\frac{x(L(t)-x)}2,
	\]
	故
	\[
		v_t(y)=p(t)+(q(t)-p(t))y+\frac{L(t)^2}2y(1-y).
	\]
	变换后的方程为 \(-L(t)^{-2}v_t''=1\)，端点值为 \(p(t),q(t)\)。密度由 \(\mathrm dx=L(t)\mathrm dy\) 变化，因此酉识别为
	\(\widehat u_t=L(t)^{1/2}v_t\)，其方程为
	\(-L(t)^{-2}\widehat u_t''=L(t)^{1/2}\)，端点值为 \(L(t)^{1/2}p(t)\)、\(L(t)^{1/2}q(t)\)。零 Dirichlet 域始终是固定的 \(H^2(0,1)\cap H^1_0(0,1)\)。

	在固定 \(y\) 处微分得到
	\[
		\dot v_t=\dot p+ (\dot q-\dot p)y+L\dot L\,y(1-y),
		\qquad -\dot v_t''=2L\dot L,
	\]
	并且 \(\dot v_t(0)=\dot p\)、\(\dot v_t(1)=\dot q\)。在可同时比较的内部点，链式法则给出
	\[
		\dot v_t(y)=(\partial_tu_t)(L(t)y)
		+u_t'(L(t)y)\dot L(t)y.
	\]
	第二项记录坐标点随区间伸缩而移动；忽略它会把变化区域的方程导数计算错。酉识别还需再加密度导数：
	\(\partial_t\widehat u_t=L^{1/2}\dot v_t+\dot L\,v_t/(2L^{1/2})\)。
\end{solution}

\begin{exercise}\label{g3:ex:coupled-spectral-crossing}
	固定 \(0<\varepsilon<1\)，考虑实对称矩阵道路
	\[
		A(t)=\begin{pmatrix}t&\varepsilon\\\varepsilon&1\end{pmatrix},
		\qquad -1\le t\le1.
	\]
	求全部零特征值时刻、相应核和交叉形式，并用两种方法计算谱流。说明耦合项如何改变交叉位置。再计算 \(B(t)=\operatorname{diag}(t^2,1)\) 的谱流，并解释为何不能对它套用非退化交叉公式。
\end{exercise}
\begin{solution}
	\(\det A(t)=t-\varepsilon^2\)，故唯一交叉在 \(t_0=\varepsilon^2\)。其核由 \((1,-\varepsilon)\) 张成。对单位核向量
	\(e=(1,-\varepsilon)/\sqrt{1+\varepsilon^2}\)，有
	\[
		\Gamma_{t_0}(e,e)=\langle\operatorname{diag}(1,0)e,e\rangle
		=\frac1{1+\varepsilon^2}>0.
	\]
	所以交叉非退化、签名为一，谱流为一。也可直接计算端点惯性：\(A(-1)\) 的行列式为负，故有一个负特征值；\(A(1)\) 的两个特征值为 \(1\pm\varepsilon>0\)，故没有负特征值。于是定义给出 \(1-0=1\)。当 \(\varepsilon=0\) 时交叉位于零，耦合将它移至 \(\varepsilon^2\)，并未改变净穿越数。

	\(B(-1)=B(1)=I\)，所以 \(\operatorname{sf}(B)=0\)。在零点的核由 \((1,0)\) 张成，但 \(\dot B(0)=0\)，交叉形式退化。特征值 \(t^2\) 仅接触零后返回正侧；它不能用一阶非退化交叉的签名公式来描述，而有限维的端点定义仍然适用。
\end{solution}
